\chapter{多机任务分配与路径规划}
\label{ch:mr-mapf}

% ════════════════════════════════════════════════════════════════
%  全吸收章：重渲染 Tutorial 笔记
%   - 04_移动机器人规控/50_多机器人协作/40_任务分配与路径规划.md
%     （任务分配：匈牙利/拍卖/SGA/CBBA；MAPF：CBS/PIBT/LaCAM；子模；ECBS/LNS；CCBS/ADG；MAPD/TAPF；学习式）
%  记号归一：R∈SO(3)、T∈SE(3)（preamble 宏 \SO \SE \so \se）；分配/搜索/图论自然记号保留。
%  源由 AI 生成，作者归属/年份/卷期/数值/前沿论文编号存疑处就地 \pz/\rebuilt，并入 claims 台账。
%  教学层遵循 v5.0 理论教学规范（动机→反面→历史→理论→陷阱→练习 + 误解汇总/小结/速查/故障排查）。
% ════════════════════════════════════════════════════════════════

\rebuilt{本章重渲染自 Tutorial 笔记《任务分配与路径规划》，其源稿由 AI 撰写，论文归属、年份、卷期、实验数值与前沿工作编号可能有误。本章忠实重述、不擅自“修正”，逐处 \texttt{\textbackslash cite} 标源；存疑断言已登记 claims 台账，待联网核对后二次校验。}

\begin{practice}[前置自测]
开始之前，先试答下列问题。若已能流畅作答，可只速读本章表格与速查卡；若卡住（答不出 $\ge 2$ 题），标 (1)(2) 的回共识/图论章（\cref{ch:mr-consensus}），标 (3)(4) 的回 A\* 与组合优化基础，再来读本章。本章会大量复用这些前置，但不重新教它们。
\begin{enumerate}
  \item 什么叫\textbf{双随机矩阵}？为什么离散共识 $\mathbf{x}(k{+}1)=\mathbf{W}\mathbf{x}(k)$ 要求 $\mathbf{W}$ 双随机才收敛到正确平均？本章 CBBA 的“冲突裁决”会用到共识的哪个性质——是收敛到平均，还是别的？
  \item 通信图的\textbf{直径}（diameter）是什么？一条信息从图一端传到另一端最少需要几轮邻居转发？这个数会直接出现在 CBBA 的收敛轮数里。
  \item A\* 搜索里\textbf{可采纳启发函数}（admissible heuristic）的定义是什么？为什么可采纳能保证 A\* 找到最优解？本章 CBS 的“低层搜索”就是一堆带约束的单体 A\*。
  \item 什么是\textbf{指派问题}（assignment problem）？给定一个 $n\times n$ 代价矩阵，把它求解到最优需要多少计算量（用大 $O$）？你听说过匈牙利算法吗？
  \item 把 $N$ 个机器人的路径规划拼成一个“联合智能体”在 $N$ 倍维度的空间里搜索，搜索空间会怎样随 $N$ 增长？这正是 MAPF 要回避的爆炸。
\end{enumerate}
\noindent\emph{答案线索}：1$\to$\cref{sec:mapf-cbba}（CBBA 的最大值共识）；2$\to$\cref{sec:mapf-cbba}（收敛随直径）；3$\to$\cref{sec:mapf-cbs}（CBS 低层时空 A\*）；4$\to$\cref{sec:mapf-assign}（LAP 与匈牙利）；5$\to$\cref{sec:mapf-def}（联合 A\* 的 $b^k$ 爆炸）。
\end{practice}

\section*{本章目标}
学完本章，你应当能够：
\begin{enumerate}
  \item \textbf{建模并实现任务分配的基线}：把多机任务分配建模为指派问题，用 Gerkey--Matari\'c 分类法（ST/MT、SR/MR、IA/TA）定位问题、说明 ST-SR-IA 为何恰可多项式最优求解；实现匈牙利（最优、$O(n^3)$）与拍卖/贪婪（次优、可分布式）两种基线，量化二者的最优性差距；
  \item \textbf{推导并实现 CBBA}：讲清两阶段结构（本地 bundle 构建 $+$ 基于共识的冲突消解），说明\textbf{边际收益递减}（DMG）为何保证它等价于集中式顺序贪婪、从而拿到 $50\%$ 最优界，并解释收敛轮数为何正比于通信图直径；
  \item \textbf{形式化 MAPF}：写出顶点冲突与边冲突的定义，区分 sum-of-costs 与 makespan，说明最优求解为何 NP-难，以及“耦合（联合 A\*）vs 解耦（单体规划）”两端的代价；
  \item \textbf{推导并实现 CBS}：讲清两层结构（高层约束树 $+$ 低层单体 A\*）、冲突如何分裂成两个加约束的子节点、为何“两子节点都展开”能保证最优，并对比优先级规划的取舍；
  \item \textbf{解释可扩展 MAPF 的工程路线}：优先级规划 $+$ 预留表、PIBT 的优先级继承与回溯、LaCAM 的惰性后继生成，说明它们用什么换取“上千智能体亚秒求解”，以及终身 MAPF 与一次性 MAPF 的区别；
  \item \textbf{统一到子模视角}：说明 $50\%$ 与 $1-1/e$ 这两个贪婪界背后是同一个子模性，把分配与路径接成一个系统，并辨析离散 MAPF 与连续编队的边界；
  \item \textbf{掌握质量$\times$规模谱}：讲清有界次优（ECBS/EECBS，解 $\le w$ 倍最优）与 anytime（MAPF-LNS，先可行再逼近最优），把全章方法收进“最优性$\times$完备性$\times$规模”总表；
  \item \textbf{把离散计划接到真机执行}：讲清连续时间 MAPF（CCBS/SIPP）、动作依赖图（ADG）与 $k$-robust，以及终身取送货 MAPD 与联合目标分配-路径 TAPF/CBM，并能为一个真实问题按“规模$\times$最优性$\times$是否终身$\times$是否上真机”选对方法。
\end{enumerate}

\subsection*{本章知识导航}
\textbf{一句话定位}：本章把“多机协调”从连续控制层往上推一层到\textbf{离散决策层}——先决定“谁去做哪件任务”（任务分配），再决定“每个机器人走哪条路、如何不在共享空间里相撞”（路径规划）。两件事都在通信图/环境图上进行，而任务分配的分布式版本（CBBA）直接调用\cref{ch:mr-consensus}的共识。

\begin{center}
\small
\begin{tikzpicture}[
  node distance=5mm and 7mm, >=Stealth,
  blk/.style={draw, rounded corners, align=center, font=\small, inner sep=3pt, fill=main!6, draw=main!60!black},
  grp/.style={draw, rounded corners, align=center, font=\small, inner sep=3pt, fill=second!6, draw=second!60!black},
  fin/.style={draw, rounded corners, align=center, font=\small, inner sep=3pt, fill=third!8, draw=third!60!black},
  ln/.style={->, thick, gray!60!black}]
\node[blk] (assign) {任务分配·集中\\ 匈牙利/拍卖/SGA};
\node[blk, right=of assign] (cbba) {任务分配·分布\\ CBBA $=$ 拍卖$+$共识};
\node[blk, right=of cbba] (mapf) {MAPF 形式化\\ 冲突/目标/NP-难};
\node[grp, below=9mm of assign] (sub) {子模视角\\ $50\%$ / $1{-}1/e$};
\node[blk, below=9mm of cbba] (cbs) {CBS 最优\\ 约束树$+$低层 A\*};
\node[grp, below=9mm of mapf] (scal) {可扩展\\ PP/PIBT/LaCAM};
\node[fin, below=18mm of cbba] (sel) {质量$\times$规模谱\\ ECBS/LNS、CCBS/ADG、选型};
\draw[ln] (assign)--(cbba); \draw[ln] (cbba)--(mapf);
\draw[ln] (assign.south)--(sub.north);
\draw[ln] (mapf.south)--(scal.north);
\draw[ln] (cbba.south)--(cbs.north);
\draw[ln] (cbs.south)--(sel.north);
\draw[ln] (scal.south)|-(sel.east);
\draw[ln] (sub.south)|-(sel.west);
\end{tikzpicture}
\end{center}

\noindent\textbf{推荐路径}：\cref{sec:mapf-assign}（分配基线，必读）$\to$\cref{sec:mapf-cbba}（CBBA，分布式方向必读、需动手实现）$\to$\cref{sec:mapf-def}（MAPF 语言，必读）$\to$\cref{sec:mapf-cbs}（CBS，必读且需搭一遍）$\to$\cref{sec:mapf-scalable}（PIBT/LaCAM，做仓库/蜂群必读）$\to$\cref{sec:mapf-submod}（分配与路径如何组合、离散与连续边界）$\to$\cref{sec:mapf-bounded}（有界次优$+$anytime，做大规模必读）$\to$\cref{sec:mapf-exec}（接真机鲁棒执行）$\to$\cref{sec:mapf-landscape}（学习式$+$总表$+$选型框架，当索引常查）。\cref{sec:mapf-cbba,sec:mapf-cbs,sec:mapf-bounded,sec:mapf-exec} 是四个硬核小节，值得花最多时间。

\subsection*{前置知识桥接}
本章紧接共识/图论，这里把最关键的前置点重新激活——你不必翻回去就能跟上：
\begin{itemize}
  \item \textbf{通信图与直径}（\cref{ch:mr-consensus}）。用图 $\mathcal{G}=(\mathcal{V},\mathcal{E})$ 建模“谁能和谁通信”，定义 Laplacian $\mathbf{L}$、代数连通性 $\lambda_2$。\cref{sec:mapf-cbba} 的 CBBA 在这张图上运行——其收敛轮数正比于图的\textbf{直径}（信息从一端传到另一端要几跳），而共识章告诉我们“图上迭代”的速度由谱（$\lambda_2$）决定。注意：任务分配的“图”和路径规划的“图”是两张不同的图——前者是机器人间的\textbf{通信图}，后者是环境的\textbf{栅格/路网图}。
  \item \textbf{共识作为“分布式裁决”的原子操作}（\cref{ch:mr-consensus}）。共识能让各 agent 只靠邻居通信收敛到一致。CBBA 的冲突消解阶段本质上是\textbf{带时间戳的最大值共识}：每个机器人维护“每个任务当前已知的最高出价”，通过邻居交换不断刷新，最终全网对“每个任务归谁”达成一致。共识在这里换层皮再次出场。
  \item \textbf{A\* 搜索}。单体最短路用 A\*：维护 OPEN 表（按 $f=g+h$ 排序）、CLOSED 表，可采纳启发 $h$ 保证最优。\cref{sec:mapf-cbs} 的 CBS “低层”就是给单个机器人跑 A\*，只是多了一组“某时刻不能在某顶点/某边”的时空约束；\cref{sec:mapf-scalable} 的优先级规划则是在“已被占用的时空”上跑 A\*。
\end{itemize}

\subsection*{如果跳过本章会怎样}
\begin{itemize}
  \item \textbf{贪婪分配看着合理，总代价却差一截}：给 $5$ 台机器人分 $5$ 个任务，图省事用“每台各自挑最近任务”。结果两台抢同一个近任务、剩下的远任务没人去，总行程比最优长出约 $40\%$。\cref{sec:mapf-assign} 会告诉你：这是贪婪/拍卖的固有次优——指派问题有最优的匈牙利（$O(n^3)$），而当你必须分布式时 CBBA 只能保证拿到最优的一半（$50\%$ 界），这是“分布式 $+$ 可扩展”必须付的税。
  \item \textbf{多机路径各自规划，跑起来全堵死}：让每台机器人独立用 A\* 规划最短路再一起跑。结果狭窄走廊里两台正面顶牛（边冲突），交叉路口同时进入同一格（顶点冲突），死锁几乎必然。\cref{sec:mapf-def,sec:mapf-cbs} 会告诉你：多机路径规划不是 $N$ 个独立单体规划——必须显式处理顶点冲突与边冲突，要么用 CBS 求最优，要么用优先级规划/PIBT 这类有协调机制的方法。
\end{itemize}

\begin{table}[H]\centering\small
\caption{本章阅读方式建议}
\label{tab:mapf-reading}
\begin{tabular}{lll}
\toprule
阅读方式 & 时间 & 适合谁 \\
\midrule
精读（推完 CBBA 的 DMG/$50\%$ 界、搭一遍 CBS、实现 PP$+$LNS$+$ADG） & 22--30 小时 & 第一次系统学多机分配与 MAPF \\
速读（读懂各算法流程与适用边界，跳过证明细节） & 7--9 小时 & 已有背景，查漏补缺 \\
速查（分配对比表 $+$ CBBA 两阶段 $+$ CBS 流程 $+$ 质量$\times$规模谱 $+$ 故障排查） & 50--60 分钟 & 实现调度/MAPF 时回查 \\
\bottomrule
\end{tabular}
\end{table}

\begin{note}
\textbf{本章记号与约定.}\;
旋转矩阵记 $\mathbf{R}\in\SO(3)$、位姿 $\mathbf{T}\in\SE(3)$（preamble 宏 $\SO,\SE,\so,\se$，与全书一致，\nameref{ch:notation}）；分配/搜索/图论领域的自然记号保留。
机器人（agent）数 $N$、任务数 $M$，下标 $i,j$。代价矩阵 $\mathbf{C}=[c_{ij}]$、效用矩阵 $\mathbf{U}=[u_{ij}]$；指派变量 $x_{ij}\in\{0,1\}$。
通信图 $\mathcal{G}_{\mathrm{c}}$、环境/栅格图 $\mathcal{G}=(V,E)$（顶点是可站位置，边是可走的相邻移动）；通信图直径记 $D$。
MAPF 中机器人 $i$ 的起点 $s_i\in V$、目标 $g_i\in V$、路径 $\pi_i=(\pi_i(0),\pi_i(1),\dots)$；$\mathrm{cost}(\pi_i)$ 为到达并停在目标所用时间步数。
\textbf{字母隔离声明}：本章 $\mathbf{J}$ 不指雅可比（流形雅可比本章不出现），$w$ 在 \cref{sec:mapf-bounded} 起指次优因子（$w\ge 1$）、其余处指权重；$\mathbf{p}_i$ 在分配语境指 CBBA 的 path、在路径语境指路径 $\pi_i$（已分别记号）；$\rho$ 不在本章出现，$\varepsilon$ 指拍卖松弛量，均仅本章局部生效。
\end{note}

% ════════════════════════════════════════════════════════════════
\section{任务分配问题与拍卖机制}
\label{sec:mapf-assign}

\paragraph{动机：谁去做哪件任务。}
设想三台巡检机器人 $\{r_1,r_2,r_3\}$ 与三个待巡检点 $\{t_1,t_2,t_3\}$，每台去每个点都有一个“代价”（行驶距离或耗时），写成代价矩阵 $\mathbf{C}$，其中 $c_{ij}$ 是机器人 $i$ 去任务 $j$ 的代价。要给每台各派一个任务（一对一），让\textbf{总代价最小}。三个任务三种排法只有 $3!=6$ 种，枚举即可；但 $N$ 个任务有 $N!$ 种分法，$N=15$ 时已超过一万亿种。更现实的是，任务分配几乎从不孤立出现——它是后续运动的入口：仓库里“哪台 AGV 去取哪个货架”、灾后搜救里“哪架无人机扫哪个区”、多臂装配里“哪只手装哪个件”。分配得好不好，直接决定系统效率上限。

\paragraph{反面：两种朴素办法都卡。}
其一，\textbf{枚举所有分法}。理论最优，但 $N!$ 的复杂度让它在 $N>12$ 左右彻底失效——这是组合爆炸，和联合优化的 $O(N^3)$ 爆炸是同一类敌人（只是更猛）。其二，\textbf{贪婪：每台挑自己最便宜的任务}。快、可分布式，但会“撞车”——两台都觉得 $t_1$ 最近、抢同一个，而没人去偏远的 $t_3$。问题的关键在于：任务分配有\textbf{全局耦合}——某台选了某任务会改变其他机器人的最优选择。朴素贪婪只看局部、忽略耦合，故次优；枚举看全局但代价爆炸。我们需要既考虑全局耦合、又多项式可解的方法——对最常见的“一对一、一次性”分配，它存在（匈牙利）；但一旦任务可被多机协作或随时间到达，问题就变难，所以先得搞清“我的问题到底是哪一类”。

\paragraph{历史：从匈牙利算法到市场机制。}
% ⚠ 复核裁决：wrong（cite_key=kuhn1955hungarian）。下面原句把 O(n^3) 系到 Kuhn(1955)/Munkres(1957)，错误——按政策注释掉原断言、不重建，仅以 \rebuilt 标错因与正确方向。
% 原文（保留备查，已停用）：任务分配的数学源头是\textbf{指派问题}：$1955$ 年 Kuhn 提出匈牙利算法（Hungarian algorithm）给出最优解\cite{kuhn1955hungarian}，Munkres $1957$ 年完善并证明其多项式复杂度\cite{munkres1957algorithms}——故又称 Kuhn--Munkres 算法，把“最小代价完美匹配”在 $O(n^3)$ 内解决。
任务分配的数学源头是\textbf{指派问题}：$1955$ 年 Kuhn 提出匈牙利算法（Hungarian algorithm）给出最优解\cite{kuhn1955hungarian}，Munkres $1957$ 年完善并证明其为（强）多项式算法\cite{munkres1957algorithms}——故又称 Kuhn--Munkres 算法。\rebuilt{错因：原句称“Munkres 1957 证明多项式复杂度、在 $O(n^3)$ 内”，把 $O(n^3)$ 归到 1955/1957 这条线，与史实不符。正确方向：Munkres(1957) 给出的原始算法是 $O(n^4)$（已证明强多项式，非 $O(n^3)$）；$O(n^3)$ 界约 15 年后才由 Edmonds--Karp(1972) 与 Tomizawa(1971) 独立达成。引用键 \texttt{kuhn1955hungarian} 本身有效；“匈牙利算法 $O(n^3)$”作为对\emph{现代变体}的描述成立，但不可作为 Kuhn/Munkres 当年所建立结果的表述。}到了机器人领域，人们更看重\textbf{可扩展}与\textbf{分布式}，于是市场/拍卖机制兴起：把任务当作待拍卖的商品，机器人作为竞标者出价，价高（或代价低）者得。Bertsekas 在 $1980$ 年代提出的拍卖算法用“价格上涨 $+$ $\varepsilon$-互补松弛”逼近最优指派\cite{bertsekas1988auction}。$2004$ 年 Gerkey 与 Matari\'c 发表了多机任务分配（MRTA）的奠基性分类法\cite{gerkey2004taxonomy}，让大家第一次能清楚地说“我研究的是哪一类问题、它的复杂度如何”；同年 Dias 等人系统综述了市场化方法\cite{dias2006market}。这条“拍卖”的线，正是 CBBA 的直接前身——CBBA $=$ 分布式拍卖 $+$ 共识。

\subsection{线性指派问题（LAP）}
\label{sec:mapf-lap}

最基础的形式：$N$ 台机器人、$N$ 个任务，一对一分配，最小化总代价。引入指派变量 $x_{ij}\in\{0,1\}$（$x_{ij}=1$ 表示机器人 $i$ 做任务 $j$），问题写成

\begin{equation}
  \min_{\mathbf{x}}\ \sum_{i=1}^{N}\sum_{j=1}^{N} c_{ij}\,x_{ij}\quad\text{s.t.}\quad
  \sum_j x_{ij}=1\ \forall i,\quad \sum_i x_{ij}=1\ \forall j,\quad x_{ij}\in\{0,1\}.
  \label{eq:mapf-lap}
\end{equation}

两组等式约束分别表示“每台机器人恰好一个任务”和“每个任务恰好一台机器人”。

\begin{theorem}[LAP 的全幺模性与多项式可解]\label{thm:mapf-tu}
\cref{eq:mapf-lap} 是整数规划，但其约束矩阵是\textbf{全幺模}（totally unimodular）的，因此线性松弛（$x_{ij}\in[0,1]$）的最优解天然是整数。这是指派问题能多项式最优求解的根本原因；匈牙利算法利用这一结构在 $O(N^3)$ 内求得最优解\cite{kuhn1955hungarian,munkres1957algorithms}。
\end{theorem}

\noindent 若用最大化效用 $u_{ij}$ 而非最小化代价，令 $c_{ij}=-u_{ij}$ 即可。

\subsection{Gerkey--Matari\'c 分类法：先认清问题属于哪一类}
\label{sec:mapf-taxonomy}

不是所有任务分配都像 LAP 那样好解。Gerkey 与 Matari\'c 用三个二元维度刻画 MRTA\cite{gerkey2004taxonomy}：

\begin{definition}[MRTA 三维分类]\label{def:mapf-taxonomy}
\begin{itemize}
  \item \textbf{ST vs MT}（single-task vs multi-task robots）：一台机器人同时只能做一个任务，还是能并行多个；
  \item \textbf{SR vs MR}（single-robot vs multi-robot tasks）：一个任务只需一台机器人，还是需要多台协作（如多机搬一件重物）；
  \item \textbf{IA vs TA}（instantaneous vs time-extended assignment）：只做一次瞬时分配，还是要为每台机器人规划一个随时间展开的任务序列（考虑未来任务）。
\end{itemize}
\end{definition}

这三维组合出 $8$ 类问题，复杂度天差地别：

\begin{itemize}
  \item \textbf{ST-SR-IA}：恰是 \cref{eq:mapf-lap} 的 LAP，匈牙利 $O(N^3)$ 最优可解——最幸运的一类。
  \item \textbf{ST-SR-TA}：每台机器人要做一串任务（序列），与车辆路径问题（VRP）相关，\textbf{NP-难}。CBBA 处理的正是这一类（每个 agent 维护一个 bundle/序列）。\pz{存疑（suspect，cite\_key=gerkey2004taxonomy）：原文并未把 ST-SR-TA 框为“VRP”——它分类为\emph{调度}问题 $R\,\|\,\sum_j w_j C_j$（全文无 VRP/vehicle routing 字样）；NP-难在原文是经 PARTITION 归约对相关变体（ALLIANCE 效率/$R\,\|\,C_{\max}$）证明的，非直接“ST-SR-TA NP-难”。CBBA（Choi 等 2009）晚于该 2004 文，不可溯源至此引用，仅可作外部并列文献。建议：归该引用时改述为“调度问题 $R\,\|\,\sum_j w_j C_j$”，VRP/mTSP/CBBA 另引后续文献。}
  \item \textbf{任何带 MR 的类}：任务需多机协作，涉及联盟形成（coalition formation），通常 NP-难。
\end{itemize}

\begin{insight}[分类法的价值]
拿到一个分配问题，先定位它的三个维度，你就立刻知道“能不能最优解、该用什么方法”。把一个 MR 任务（需两机协作）错当成 SR 来分配，是初学者最常见的建模错误。
\end{insight}

\subsection{拍卖算法：可分布式的次优解法}
\label{sec:mapf-auction}

拍卖算法把分配看成市场：每个任务 $j$ 有当前价格 $p_j$，机器人 $i$ 对任务 $j$ 的“净收益”是 $u_{ij}-p_j$。每轮，未分配的机器人挑净收益最高的任务竞标，把价格抬高到刚好让它仍是最优、且超过次优一个 $\varepsilon$。价格单调上涨，过程在有限步收敛到一个 $\varepsilon$-近似最优指派\cite{bertsekas1988auction}：

\begin{theorem}[拍卖算法的 $\varepsilon$-最优性]\label{thm:mapf-auction}
拍卖算法的总效用与最优指派的差距不超过 $N\varepsilon$。证明从 $\varepsilon$-互补松弛条件出发（每个机器人分到的任务净收益在最优的 $\varepsilon$ 邻域内），对 $N$ 台机器人求和即得\cite{bertsekas1988auction}。
\end{theorem}

\begin{insight}[价格作为全局协调信号]
拍卖算法的妙处在于它\textbf{天然分布式}——每个机器人只需知道任务价格（全局共享的少量标量）就能独立竞标，不需要中央求解器看到整个代价矩阵。这正是它在机器人领域比匈牙利更受欢迎的原因，也是 CBBA 的直接灵感。
\end{insight}

上面是 Bertsekas 式的“单任务竞价拍卖”。机器人领域还有一类更贴近多任务场景的市场机制——\textbf{顺序单项拍卖}（Sequential Single-Item auction, SSI）\cite{koenig2006ssi}：把任务一个个拿出来拍，每轮所有机器人对\textbf{当前这件}任务报价（报价 $=$ 把该任务加进自己已有任务集后的\textbf{边际代价}），出价最低者赢得该任务、纳入自己的集合，再拍下一件。SSI 逐件拍而非一次性整体优化（比集中最优快、可分布式），但每件的报价考虑了“已接任务”的协同（比纯贪婪质量高），是仓库与多机系统里常见的实用折中。

\subsection{顺序贪婪（SGA）与“边际收益递减”}
\label{sec:mapf-sga}

当一台机器人要做多个任务（TA），它对一个任务的价值依赖于它已经接了哪些任务——比如已经要去东边了，再接一个东边的任务边际成本就低。\textbf{顺序贪婪算法}（Sequential Greedy Algorithm, SGA）每轮在所有“(机器人, 任务)对”里挑边际收益最高的一对、分配掉，重复直到无任务可分。SGA 是集中式的，但它有一个关键性质：

\begin{definition}[边际收益递减 DMG]\label{def:mapf-dmg}
称效用满足\textbf{边际收益递减}（Diminishing Marginal Gain, DMG），若一个任务的边际价值不会因为先接了别的任务而\textbf{增加}，即对路径 $\mathbf{p}\subseteq\mathbf{p}'$ 有 $c_{ij}[\mathbf{p}]\ge c_{ij}[\mathbf{p}']$。这正是组合优化里的\textbf{子模性}（见 \cref{sec:mapf-submod}）在序列上的体现。
\end{definition}

\begin{theorem}[SGA 的 $50\%$ 界]\label{thm:mapf-sga}
当效用满足 DMG（\cref{def:mapf-dmg}）时，SGA 的解保证不差于最优解的一半。这个 $50\%$ 界，以及 SGA 与 CBBA 的等价，是 \cref{sec:mapf-cbba} CBBA 拿到性能保证的全部理论依托。
\end{theorem}

把本节四种任务分配方法摆在一起，坐标是“集中还是分布 $\times$ 最优还是次优”：

\begin{table}[H]\centering\small
\caption{四种任务分配方法对比}
\label{tab:mapf-assign-compare}
\begin{tabular}{lllll}
\toprule
方法 & 集中/分布 & 最优性 & 通信/计算 & 适用 \\
\midrule
匈牙利（Kuhn--Munkres） & 集中（需全代价矩阵） & 最优（一对一） & $O(n^3)$，需全局信息 & 小规模、一对一、要最优 \\
拍卖（Bertsekas） & 可分布（只需共享价格） & $\varepsilon$-近似（差距 $\le N\varepsilon$） & 多轮竞价，通信少 & 要可分布、近最优 \\
SGA（顺序贪婪） & 集中 & $\ge\tfrac12$ 最优（DMG 下） & 每轮全局挑最优对 & 多任务、要保证、可集中 \\
SSI（顺序单项拍卖） & 可分布 & 启发式（介于贪婪与最优） & 逐件竞价，实现简单 & 多任务、分布式、实用折中 \\
\bottomrule
\end{tabular}
\end{table}

这张表是 \cref{sec:mapf-cbba} 的跳板：CBBA 要做的，正是把“SGA 的 $50\%$ 保证”和“拍卖/SSI 的可分布”合二为一——用共识让分布式的竞价结果等价于集中式 SGA。

\subsection{代码验证：匈牙利 vs 贪婪的最优性差距}
\label{sec:mapf-assign-code}

用代码把“贪婪次优、匈牙利最优”看清楚。\verb@scipy.optimize.linear_sum_assignment@ 就是匈牙利算法的实现。

\begin{codebox}[Python]{匈牙利 vs 贪婪 vs 拍卖：量化最优性差距}
import numpy as np
from scipy.optimize import linear_sum_assignment
rng = np.random.default_rng(0)

def greedy_assignment(C):
    """贪婪：每轮挑全局最小代价的(机器人,任务)对，分配掉。"""
    n = C.shape[0]; C = C.copy().astype(float)
    rows, cols = set(range(n)), set(range(n)); total, assign = 0.0, {}
    while rows:
        sub = C[np.ix_(list(rows), list(cols))]
        ri, ci = np.unravel_index(np.argmin(sub), sub.shape)
        i, j = list(rows)[ri], list(cols)[ci]
        assign[i] = j; total += C[i, j]; rows.remove(i); cols.remove(j)
    return total, assign

def hungarian_cost(C):
    r, c = linear_sum_assignment(C); return C[r, c].sum()     # 最优

for n in [3, 5, 10, 20]:
    gaps = []
    for _ in range(200):
        C = rng.uniform(1, 100, size=(n, n))
        opt = hungarian_cost(C); gre, _ = greedy_assignment(C)
        gaps.append((gre - opt) / opt * 100)                  # 贪婪比最优贵百分之几
    print(f"n={n:2d}: 贪婪平均比最优贵 {np.mean(gaps):5.1f}% , 最坏 {np.max(gaps):5.1f}%")

def auction_assignment(U, eps=1e-3):
    """拍卖算法(最大化效用)。每个机器人只看任务价格 p 竞标——可分布式的关键。"""
    n = U.shape[0]; p = np.zeros(n); owner = -np.ones(n, dtype=int)
    unassigned = list(range(n))
    while unassigned:
        i = unassigned.pop(); net = U[i] - p
        j = int(np.argmax(net)); best = net[j]
        second = np.partition(net, -2)[-2]                    # 次优净收益
        bid = p[j] + (best - second) + eps                    # 抬价：超过次优一个 eps
        if owner[j] != -1: unassigned.append(owner[j])        # 原主被挤掉，重新竞标
        owner[j] = i; p[j] = bid
    return sum(U[owner[j], j] for j in range(n)), owner
\end{codebox}

\textbf{运行要点}：你会看到贪婪平均比最优贵百分之十几到几十，且\textbf{最坏情况差距随规模拉大}——这量化了“局部贪婪 vs 全局最优”的代价。匈牙利之所以能最优，是因为它利用了 \cref{eq:mapf-lap} 的全幺模结构（\cref{thm:mapf-tu}）做全局优化；贪婪每步只看当前最小，丢掉了全局耦合。拍卖的总效用与最优只差极小一点（由 $\varepsilon$ 控制、理论上界 $N\varepsilon$，\cref{thm:mapf-auction}），却只靠“每个机器人看价格竞标”实现——没有任何一处需要看到完整效用矩阵。这个“价格作为全局协调信号、各 agent 独立竞标”的结构，就是 \cref{sec:mapf-cbba} CBBA 的骨架，只不过 CBBA 把“价格”换成了通过\textbf{共识}传播的“中标信息”。

\begin{pitfall}[概念误区：把多机协作任务（MR）当成单机任务（SR）来分配]
有个任务实际需要两台机器人协作（如合抬一件重物），却用 LAP/匈牙利做一对一分配，给它派一台机器人。后果：分配出的方案物理上不可执行，或系统在执行阶段才发现要临时拉人。根本原因：LAP 的模型假设是 ST-SR-IA；MR 任务属于联盟形成问题，数学结构完全不同（NP-难、需考虑机器人组合）。正确做法：动手前先用 \cref{def:mapf-taxonomy} 的三维给问题归类，带 MR 的问题不能用 LAP。
\end{pitfall}

\begin{pitfall}[思维陷阱：默认“最小化总代价”就是唯一目标]
不加思考地把目标设成总行程/总耗时之和（sum），用匈牙利求解。后果：得到的方案总和最小，但某一台被派了极远任务、单独干到天黑（完工时间 makespan 很大），团队整体要等它。根本原因：sum 与 makespan（最大单体完工时间）是不同目标——前者优化吞吐，后者优化“最后一个干完的时间”；匈牙利默认优化 sum。正确做法：要均衡负载/抢时间用 min-max（makespan），它一般比 min-sum 更难（常需二分 $+$ 可行性检查）。
\end{pitfall}

\begin{pitfall}[编程陷阱：代价矩阵的行列方向或非方阵处理出错]
用 \verb@linear_sum_assignment@ 时把“机器人$\times$任务”的矩阵转置了，或机器人数 $\ne$ 任务数时直接塞进去不补齐。后果：分配结果张冠李戴，或在非方阵上得到部分分配却没意识到有任务/机器人被漏掉。根本原因：\verb@linear_sum_assignment(C)@ 约定 \verb@C[i,j]@ 为第 $i$ 行（行$=$机器人）指派到第 $j$ 列（列$=$任务）的代价，返回 \verb@row_ind, col_ind@ 配对；非方阵时它做最大匹配。正确做法：统一约定“行$=$机器人、列$=$任务”并写进注释；机器人数 $\ne$ 任务数时显式补齐（用大代价的虚拟行/列）；务必用返回的 \verb@(row_ind, col_ind)@ 成对取值。
\end{pitfall}

\begin{exercise}
\begin{enumerate}
  \item \textbf{[实现]} 实现并对比匈牙利、贪婪、拍卖三种分配。在 $n\in\{5,10,20,50\}$ 上各跑 $200$ 次随机代价矩阵，画出“贪婪/拍卖相对最优的差距百分比”随 $n$ 的曲线。验证：拍卖差距由 $\varepsilon$ 控制且很小，贪婪差距随 $n$ 变大。
  \item \textbf{[建模]} 对下面三个场景用 \cref{def:mapf-taxonomy} 三维归类并说明能否用匈牙利最优求解：(a) $5$ 台扫地机各扫一个房间；(b) $3$ 台无人机轮流巡逻 $8$ 个点位；(c) $4$ 台四足合抬 $2$ 件重物（每件需 $2$ 机）。
  \item \textbf{[实现/分析]} 把目标从 min-sum 改成 min-max（makespan）：实现一个“二分 $+$ 可行性检查（用匈牙利判定）”的 min-max 求解器，对比两种目标给出的分配差异。
  \item \textbf{[推导]} 证明 \cref{thm:mapf-auction}：拍卖算法的总效用与最优指派的差距不超过 $N\varepsilon$。提示：从 $\varepsilon$-互补松弛条件出发对 $N$ 台机器人求和。
  \item \textbf{[思考，接 \cref{sec:mapf-cbba}]} 拍卖里“价格 $p_j$”是全局共享的协调信号。设想没有中央黑板维护价格、每台机器人只能和邻居通信，它们要怎样才能就“每个任务当前的最高出价是多少、归谁”达成一致？
\end{enumerate}
\end{exercise}

% ════════════════════════════════════════════════════════════════
\section{CBBA：基于共识的分布式拍卖}
\label{sec:mapf-cbba}

\paragraph{动机：没有中央黑板的拍卖。}
\cref{sec:mapf-auction} 的拍卖算法虽然“每个机器人看价格独立竞标”，但那个价格仍假设有一块全局共享的黑板。真实多机系统没有黑板——每台机器人只能和通信范围内的几个邻居说话（这正是\cref{ch:mr-consensus}的通信图）。我们想要的是：每台机器人\textbf{自己决定}接哪些任务，通过和邻居反复交换“我打算接哪些、出价多少”，最终全队达成一个\textbf{无冲突}（每个任务至多一台机器人）且\textbf{总效用不太差}的分配——全程不依赖任何中央节点。这件事难在“无冲突”是个全局性质，而每台机器人只有局部信息：两台相距很远、暂时不能直接通信的机器人，可能各自都把同一个任务加进了自己的清单。CBBA 给出的答案是：用\textbf{共识}——让“每个任务当前最高出价是谁”这个信息像共识一样在图上传播，直到全网一致。

\paragraph{反面：朴素分布式贪婪会冲突、会活锁。}
最朴素的分布式做法：每台机器人各自贪婪挑任务，然后执行。问题立刻浮现：\textbf{冲突}——机器人 $r_1$ 和 $r_3$ 可能都把 $t_5$ 收进了清单，执行时两个一起冲向 $t_5$。退一步，加个“先到先得”：但在只有邻居通信的图上，“先到”本身就模糊——$r_1$ 的宣告要好几跳才能传到 $r_3$，在这期间 $r_3$ 早已基于过时信息做了决定。更糟的是，若消解规则设计不当（允许机器人随意反悔、释放任务的顺序混乱），系统会陷入\textbf{活锁}（livelock）：任务在机器人间反复易手、永不稳定。朴素做法缺两样东西：一是一致的、带时间戳的冲突裁决规则；二是对“释放任务”的纪律约束。CBBA 正是把这两样补齐——前者用带时间戳的共识，后者用 bundle 的有序结构 $+$ DMG 条件。

\paragraph{历史：Choi、Brunet、How 的 CBBA。}
CBBA 由 Choi、Brunet 与 How 于 $2009$ 年在 \emph{IEEE Transactions on Robotics} 上提出（《Consensus-Based Decentralized Auctions for Robust Task Allocation》）\cite{choi2009cbba}。它有两个版本：CBAA 处理单任务分配（每机一个任务），CBBA 处理多任务分配（每机一串任务，即 bundle）。论文的核心贡献是：把拍卖的分布式竞标与共识的冲突消解结合，并\textbf{证明}了在 DMG 条件下 CBBA 产生的解与集中式顺序贪婪（SGA）完全相同，从而继承 SGA 的 $50\%$ 最优界，且在有限轮数内收敛到无冲突。这一组合至今仍是分布式 MRTA 的基准。

\subsection{两个数据结构：bundle 与 path}
\label{sec:mapf-cbba-struct}

每台机器人 $i$ 维护两个有序列表：

\begin{itemize}
  \item \textbf{bundle} $\mathbf{b}_i$：按\textbf{加入顺序}排列的任务列表——记录“我是按什么顺序把这些任务收进来的”；
  \item \textbf{path} $\mathbf{p}_i$：按\textbf{执行顺序}排列的同一批任务——记录“我打算按什么顺序去做”。
\end{itemize}

区别在于：加入 bundle 时，新任务可以插进 path 的任意位置（选使总收益最大的插入点），但在 bundle 里永远追加在末尾。这个区别至关重要——冲突消解时按 \textbf{bundle 顺序}释放任务，正是保证 DMG 性质成立的关键。

\subsection{阶段一：bundle 构建（本地贪婪）}
\label{sec:mapf-cbba-build}

每台机器人独立地、贪婪地往自己的 bundle 里加任务。给定当前 path $\mathbf{p}_i$，任务 $j$ 的\textbf{边际得分} $c_{ij}[\mathbf{p}_i]$ 定义为“把 $j$ 插入 $\mathbf{p}_i$ 的最佳位置所能带来的得分增量”\cite{choi2009cbba}：

\begin{equation}
  c_{ij}[\mathbf{p}_i]=\max_{k\le |\mathbf{p}_i|}\ \big[\,S_i^{\mathbf{p}_i\oplus_k\{j\}}-S_i^{\mathbf{p}_i}\,\big],
  \label{eq:mapf-cbba-marginal}
\end{equation}

其中 $S_i^{\mathbf{p}}$ 是机器人 $i$ 执行路径 $\mathbf{p}$ 的总得分（如“任务奖励之和减去行程代价”），$\oplus_k$ 表示在第 $k$ 个位置插入。机器人每轮挑边际得分最高、且高于该任务当前已知最高出价的任务加入，直到 bundle 满（达到容量 $L_t$）或没有任务能再提升得分。每次加入都记下自己对该任务的出价 $y_{ij}=c_{ij}$。

\subsection{阶段二：基于共识的冲突消解}
\label{sec:mapf-cbba-consensus}

每台机器人维护三个向量（长度 $=$ 任务数 $M$）：

\begin{itemize}
  \item \textbf{中标价} $\mathbf{y}_i\in\mathbb{R}^M$：$y_i(j)$ $=$ 机器人 $i$ 当前\textbf{已知}的任务 $j$ 的最高出价；
  \item \textbf{中标者} $\mathbf{z}_i\in\{1,\dots,N\}^M$：$z_i(j)$ $=$ 机器人 $i$ 当前认为任务 $j$ 归谁；
  \item \textbf{时间戳} $\mathbf{s}_i$：记录 $i$ 最近一次从各机器人收到信息的时间（用于裁决信息新旧）。
\end{itemize}

每个通信周期，机器人和邻居交换 $(\mathbf{y},\mathbf{z},\mathbf{s})$，然后对每个任务按一组\textbf{确定性规则}更新（论文给出完整的“action rules”表）：若邻居对任务 $j$ 的中标价更高，就更新 $y_i(j),z_i(j)$ 为邻居的值；若发现自己 bundle 里某个任务被别人以更高价抢走（被 outbid），就\textbf{把该任务连同它之后（在 bundle 里）的所有任务一起释放}，然后回到阶段一重新竞标。把论文的规则归纳一下，接收方 $i$ 收到邻居 $k$ 关于任务 $j$ 的信息后，无非做三类动作之一——\textbf{更新}（采纳 $k$ 的 $y,z$）、\textbf{重置}（把 $y_i(j)$ 清零、$z_i(j)$ 置空）、\textbf{保留}（维持自己的）。具体做哪个，由“发送方 $k$ 认为 $j$ 归谁”和“接收方 $i$ 认为 $j$ 归谁”两维共同决定：

\begin{table}[H]\centering\small
\caption{CBBA action rules 精简表（接收方 $i$ 收到邻居 $k$ 关于任务 $j$ 的信息后的动作；行 $=$ 发送方 $k$ 认为归谁，列 $=$ 接收方 $i$ 认为归谁）}
\label{tab:mapf-cbba-rules}
\begin{tabular}{lllll}
\toprule
$k$ 认为 $\backslash$ $i$ 认为 & $i$ 自己 & $k$ 自己 & 第三方 $m$ & 无人 \\
\midrule
$k$ 自己 & $y_k>y_i\Rightarrow$ 更新 & 更新 & $y_k>y_i$ 或 $k$ 更新 $\Rightarrow$ 更新 & 更新 \\
$i$ 自己 & 保留 & 重置 & $k$ 信息更新 $\Rightarrow$ 重置 & 保留 \\
第三方 $m$ & $y_k>y_i\Rightarrow$ 更新 & 更新 & $k$ 更新且 $y_k>y_i\Rightarrow$ 更新 & 更新 \\
无人 & 保留 & 更新 & $k$ 信息更新 $\Rightarrow$ 重置 & 保留 \\
\bottomrule
\end{tabular}
\end{table}

读这张表的钥匙是：\textbf{发送方对“自己中标的任务”最有发言权}（第一行几乎都更新），而\textbf{“重置”专门处理信念矛盾}——当 $k$ 说“这任务归你 $i$”、$i$ 却以为归别人时，说明 $i$ 的信念已过时，清零重判最稳妥。本节代码用同步精简版（直接取“出价更高者胜、等价按编号裁决”），它抓住了上表最主导的“更高出价者胜”情形；真实异步 CBBA 用完整的时间戳 $+$ 上表来应对消息乱序、丢包与信息不一致。

\begin{insight}[CBBA 的冲突消解就是最大值共识]
这一步的本质是\cref{ch:mr-consensus}的\textbf{最大值共识}：把“每个任务的最高出价” $\mathbf{y}$ 当作要共识的量，通过邻居间反复取最大值（带时间戳裁决），最终全网对 $\mathbf{y},\mathbf{z}$ 达成一致——而 $\mathbf{y},\mathbf{z}$ 一致就意味着所有机器人对“每个任务归谁”看法一致，即无冲突。共识章“图谱决定迭代速度”的结论，在这里体现为：\textbf{收敛轮数正比于通信图直径}。
\end{insight}

\subsection{为什么“按 bundle 顺序释放”？DMG 与 $50\%$ 界}
\label{sec:mapf-cbba-bound}

这是 CBBA 最精妙之处。当机器人被 outbid 而要释放任务 $j$ 时，它必须把 $j$ \textbf{以及 bundle 里排在 $j$ 之后的所有任务}一并释放，哪怕那些后续任务没被任何人抢。

\begin{insight}[释放纪律的理由]
后加入任务的边际得分，是在“$j$ 已在 path 中”的前提下算出来的（见 \cref{eq:mapf-cbba-marginal}）；一旦 $j$ 被拿走，这些后续任务的得分就不再有效。只有按加入顺序释放，才能保证每个保留下来的任务的出价始终反映其真实边际价值。
\end{insight}

正是这条纪律，让 CBBA 等价于集中式 SGA。结合 \cref{thm:mapf-sga}：

\begin{theorem}[CBBA $=$ SGA 与 $50\%$ 界]\label{thm:mapf-cbba}
在效用满足 DMG（\cref{def:mapf-dmg}）时，CBBA 产生的解与集中式顺序贪婪 SGA 完全相同，从而保证不差于最优的一半\cite{choi2009cbba}：
\begin{equation}
  S(\mathrm{CBBA})=S(\mathrm{SGA})\ \ge\ \tfrac12\,S(\mathrm{OPT}).
  \label{eq:mapf-cbba-half}
\end{equation}
\end{theorem}

\begin{proof}[证明思路]\rebuilt{（据 \cite{choi2009cbba}）}
直觉：贪婪每步拿走的边际收益，至少是它“挡掉”的最优分配中对应收益的一半（因为 DMG 保证后者的边际价值只会更低）。若效用不满足 DMG，可对得分做单调化变换（warping）强制其满足，但会牺牲一些最优性。
\end{proof}

\subsection{收敛性：正比于通信图直径}
\label{sec:mapf-cbba-converge}

CBBA 的收敛性可这样理解：对单个任务，全网就“它归谁”达成一致所需轮数不超过通信图直径 $D$（信息从中标者传到最远机器人要 $D$ 跳）；$M$ 个任务的冲突在最坏情况下需要 $O(N_{\min}D)$ 轮（$N_{\min}=\min(M,NL_t)$ 为可分配任务上限）\cite{choi2009cbba}。这把共识章“图谱决定迭代速度”的结论具体化为一句可操作的话：\textbf{通信图越“扁”（直径越小），CBBA 收敛越快}。

\paragraph{走一遍最小实例。}
$2$ 台机器人、$2$ 个任务，每机容量 $1$，通信图是一条链 $R_1\!-\!R_2$（直径 $1$）。出价矩阵（机器人对任务的得分）为 $R_1:[9,8]$、$R_2:[7,6]$——两台都最想要任务 $1$。跟着 CBBA 跑：

\begin{table}[H]\centering\small
\caption{CBBA 在 $2$ 机 $2$ 任务链上的逐轮演化}
\label{tab:mapf-cbba-trace}
\begin{tabular}{lllll}
\toprule
轮次 & 阶段 & $R_1$ 的 $(\mathbf{y},\mathbf{z})$ & $R_2$ 的 $(\mathbf{y},\mathbf{z})$ & 说明 \\
\midrule
$0$ & bundle 构建 & $y{=}[9,-],z{=}[1,-]$ & $y{=}[7,-],z{=}[2,-]$ & 两台各自拍下任务 $1$ \\
$1$ & 共识（交换） & $y{=}[9,-],z{=}[1,-]$ & $y{=}[9,-],z{=}[1,-]$ & $R_2$ 见 $9>7$，更新：任务 $1$ 归 $R_1$ \\
$1$ & 释放 $+$ 重建 & $z{=}[1,-]$ 不变 & 释放任务 $1$，改拍任务 $2$ & $R_2$ 被 outbid，转拿次优任务 $2$ \\
$2$ & 共识 & $y{=}[9,6],z{=}[1,2]$ & $y{=}[9,6],z{=}[1,2]$ & 全网一致，无冲突 \\
\bottomrule
\end{tabular}
\end{table}

结果与集中式 SGA 完全一致（SGA 也是先把全局最高的 $9$ 即“$R_1$-任务$1$”定下，再定剩下的“$R_2$-任务$2$”）——这正是“CBBA $=$ SGA”的微观写照，且收敛只用了与直径同量级的轮数。

\subsection{代码验证：CBBA 无冲突 $+$ $50\%$ 界 $+$ 收敛随直径}
\label{sec:mapf-cbba-code}

下面实现一个精简版 CBBA，在一张通信图上运行，验证三件事：最终无冲突、总得分不差于最优一半、收敛轮数随图直径增长。

\begin{codebox}[Python]{精简 CBBA：在完全图/环/链三种拓扑上验证}
import numpy as np, networkx as nx
rng = np.random.default_rng(2)

def cbba(reward, comm_graph, capacity):
    """精简 CBBA。reward[i,j]: 机器人 i 完成任务 j 的得分(此处用 DMG 的简单可加得分)。
    返回:每机的任务集合、总得分、收敛所用通信轮数。"""
    N, M = reward.shape
    y = np.zeros((N, M)); z = -np.ones((N, M), dtype=int)
    bundle = [[] for _ in range(N)]

    def build_bundle(i):                                   # 阶段一:贪婪加任务
        while len(bundle[i]) < capacity:
            cand = [j for j in range(M) if j not in bundle[i] and reward[i, j] > y[i, j]]
            if not cand: break
            j = max(cand, key=lambda j: reward[i, j])      # 可加得分下边际即 reward
            bundle[i].append(j); y[i, j] = reward[i, j]; z[i, j] = i

    for i in range(N): build_bundle(i)
    rounds = 0; changed = True
    while changed:
        changed = False; rounds += 1; new_y, new_z = y.copy(), z.copy()
        for i in range(N):                                 # 阶段二:带编号裁决的最大值共识
            for k in comm_graph.neighbors(i):
                for j in range(M):
                    if (new_y[i, j], -new_z[i, j]) < (y[k, j], -z[k, j]):
                        new_y[i, j] = y[k, j]; new_z[i, j] = z[k, j]
        if not np.array_equal(new_y, y) or not np.array_equal(new_z, z): changed = True
        for i in range(N):                                 # 释放被抢任务(连同 bundle 中其后)
            release_from = None
            for idx, j in enumerate(bundle[i]):
                if new_z[i, j] != i: release_from = idx; break
            if release_from is not None:
                for j in bundle[i][release_from:]:
                    if new_z[i, j] == i: new_y[i, j] = 0; new_z[i, j] = -1
                bundle[i] = bundle[i][:release_from]; changed = True
        y, z = new_y, new_z
        for i in range(N):
            before = len(bundle[i]); build_bundle(i)
            if len(bundle[i]) != before: changed = True
        if rounds > 200: break
    assign = {i: list(bundle[i]) for i in range(N)}
    total = sum(reward[i, j] for i in range(N) for j in bundle[i])
    return assign, total, rounds

N, M, cap = 5, 8, 2
reward = rng.uniform(1, 50, size=(N, M))
upper = sum(reward[:, j].max() for j in range(M))          # 松弛上界
for diam_graph, name in [(nx.complete_graph(N), "完全图(直径1)"),
                         (nx.cycle_graph(N), "环(直径2)"), (nx.path_graph(N), "链(直径4)")]:
    assign, total, rounds = cbba(reward, diam_graph, cap)
    all_tasks = [j for i in assign for j in assign[i]]
    conflict = len(all_tasks) != len(set(all_tasks))
    print(f"{name}: 收敛 {rounds:2d} 轮 | 直径 {nx.diameter(diam_graph)} "
          f"| 无冲突={not conflict} | 总得分 {total:.1f} (松弛上界 {upper:.1f})")
\end{codebox}

\textbf{运行要点}：三种拓扑都收敛到\textbf{无冲突}分配（每个任务至多一台机器人），且总得分相对松弛上界处于合理区间（可加得分下贪婪通常远好于 $50\%$ 的最坏界，但理论保证就是 $50\%$，见 \cref{eq:mapf-cbba-half}）。关键观察是\textbf{收敛轮数随直径单调增长}——完全图（直径 $1$）几轮搞定，链（直径 $4$）需要更多轮。这就把共识章“图谱/直径决定迭代速度”的抽象结论，变成了你能在屏幕上读出的数字。

\begin{pitfall}[概念误区：以为 CBBA 给的是最优分配]
把 CBBA 当成“分布式的匈牙利”，期望它给出和集中最优一样的解，在对比实验里发现 CBBA 总效用低于匈牙利、误以为实现有 bug。根本原因：CBBA 等价于集中式\textbf{贪婪}（SGA），不是最优；理论保证是“不差于最优的一半”（DMG 条件下，\cref{thm:mapf-cbba}），这是分布式 $+$ 可扩展必须付出的代价。正确做法：把 CBBA 和“集中贪婪 SGA”对比（应当几乎一致），而不是和匈牙利对比；$50\%$ 是\textbf{最坏界}，典型实例往往好得多。
\end{pitfall}

\begin{pitfall}[思维陷阱：释放任务时不按 bundle 顺序、只释放被抢的那一个]
机器人被 outbid 某任务 $j$ 时，只把 $j$ 释放掉，保留 bundle 里 $j$ 之后的任务。后果：保留任务的出价不再有效（它们是在“$j$ 还在”的前提下算的边际价值），导致分配次优甚至持续抖动、活锁。根本原因：后加入任务的边际得分依赖于先加入的任务（\cref{eq:mapf-cbba-marginal}）；$j$ 一走，其后任务的得分前提就塌了。正确做法：严格按 bundle（加入顺序）释放——从被抢任务起，把它及其后所有任务全部释放再重新竞标（代码里用切片 \verb@bundle[i][release_from:]@，别用 path 顺序）。
\end{pitfall}

\begin{pitfall}[编程陷阱：得分函数不满足 DMG 却直接套用 $50\%$ 界]
用了一个边际得分会\textbf{随已接任务增多而上升}的效用（如“任务越多协同奖励越高”），仍指望 CBBA 收敛且有 $50\%$ 保证。后果：CBBA 可能不收敛（反复加任务），$50\%$ 界完全失效。根本原因：$50\%$ 界与“CBBA $=$ SGA”的等价都以 DMG 为前提；非 DMG（超可加、协同增益）破坏了贪婪的次模性。正确做法：先验证得分函数是否 DMG；常见的“奖励固定 $-$ 行程代价”型通常 DMG；若有协同增益，要么改用专门方法，要么对得分做单调化 warping（牺牲部分最优性）。
\end{pitfall}

\begin{exercise}
\begin{enumerate}
  \item \textbf{[实现]} 在上面 CBBA 代码基础上，系统扫描通信图直径（完全图、环、链、星）与机器人数 $N$，画出“收敛轮数 vs 直径”的关系图，验证轮数大致正比于直径，并和共识章“收敛随图谱”的结论对照。
  \item \textbf{[实现/验证]} 构造一个集中式 SGA（每轮在所有未分配的(机器人,任务)对里挑全局边际最高的一对），在同一批随机实例上验证 \textbf{CBBA 的解与 SGA 完全相同}（坐实 \cref{thm:mapf-cbba}）。
  \item \textbf{[实现/反例]} 故意构造一个\textbf{违反 DMG} 的得分函数（机器人接的任务越多、每个任务的得分越高），在 CBBA 里运行：它是否收敛？总得分相对最优差多少？再对得分做单调化 warping 后重跑对比。
  \item \textbf{[推导]} 证明 \cref{eq:mapf-cbba-half}：在 DMG 条件下 SGA 的总效用不小于最优的一半。提示：对贪婪每步选中的对，论证它“挤掉”的最优分配里的对的边际收益不超过它本身（用 DMG）。
  \item \textbf{[设计，接 \cref{sec:mapf-cbs}]} CBBA 解决“谁去做哪个任务”，但没说“怎么走过去不撞上”。设计把 CBBA 输出（每台机器人一串任务及其顺序）交给 MAPF 求解器的接口：CBBA 该输出什么、MAPF 需要什么输入、若 MAPF 发现某分配根本走不通（路径死锁）该如何反馈给 CBBA 重分配。
\end{enumerate}
\end{exercise}

% ════════════════════════════════════════════════════════════════
\section{多机路径规划（MAPF）与最优性}
\label{sec:mapf-def}

\paragraph{动机：从“谁去做”到“怎样走不撞”。}
\cref{sec:mapf-assign,sec:mapf-cbba} 决定了“谁去做哪个任务”。但任务往往在不同地点，机器人要\textbf{走过去}——一群机器人在共享空间里移动，如何保证谁也不撞谁？如果只有一台机器人，这就是单体最短路，A\* 即可。但\textbf{一群}机器人同时在动，问题质变了：它们共享同一片空间，两台不能同时占据同一个格子，也不能在同一时刻交换位置（否则就是穿模或正面相撞）。这就是 MAPF（Multi-Agent Path Finding）——给定一张图、若干机器人各自的起点与目标，为\textbf{所有}机器人同时找到无碰撞的路径。它是仓库机器人（亚马逊 Kiva 系统调度上万台 AGV）、自动化港口、游戏寻路、无人机蜂群的核心问题。MAPF 看似只是“单体 A\* 做 $N$ 次”，但正是“无碰撞”这个全局耦合，让它从多项式问题跃升为 NP-难。

\paragraph{反面：两种朴素极端正好是这个领域的两端。}
其一，\textbf{各自独立 A\*}：每台机器人无视他人、各规划自己的最短路，然后一起跑。问题：必然碰撞，且谁也不知道该让谁，系统死锁。这是\textbf{解耦}（decoupled）的极端：快，但既不完备也不安全。其二，\textbf{联合智能体 A\*}：把 $N$ 台机器人看成一个“超级智能体”，它的状态是所有机器人位置的组合 $(v_1,v_2,\dots,v_N)$，在这个 $N$ 维联合空间里跑 A\*。问题：状态空间随 $N$ \textbf{指数爆炸}——每台机器人每步有约 $5$ 个动作（上下左右 $+$ 等待），联合分支因子就是 $5^N$，$N=10$ 时单步就近千万种组合。这是\textbf{耦合}（coupled）的极端：最优、完备，但计算上不可行。MAPF 的全部技术张力就在这两端之间：解耦快但会撞，耦合对但会爆。CBS（\cref{sec:mapf-cbs}）的精妙在于——它\textbf{搜索时是单体的（像解耦），但通过约束逐步消解冲突来保证最优（像耦合）}。

\paragraph{历史：从拼图问题到现代 MAPF。}
MAPF 的理论源头是组合数学里的“华容道/拼图”问题：$1984$ 年 Kornhauser 等人研究了图上的“卵石移动”（pebble motion），给出了可解性判定\cite{kornhauser1984pebble}。机器人领域真正把它做成系统问题，是亚马逊收购的 Kiva Systems——Wurman 等人 $2008$ 年描述了如何协调成百上千台仓库机器人\cite{wurman2008coordinating}。理论上，Yu 与 LaValle 在 $2013$ 年证明了 MAPF 在 sum-of-costs 与 makespan 目标下的\textbf{最优求解都是 NP-难}\cite{yu2013structure}。$2019$ 年 Stern 等人发表了 MAPF 的标准术语与基准综述\cite{stern2019mapf}，统一了顶点/边冲突的定义、目标函数、地图集合——此后 MAPF 研究有了公共的比较基准。

\subsection{MAPF 的形式化定义}
\label{sec:mapf-def-formal}

\begin{definition}[MAPF 实例与解]\label{def:mapf}
给定无向图 $\mathcal{G}=(V,E)$（顶点是可站位置，边是可走的相邻移动）和 $k$ 台机器人，机器人 $i$ 有起点 $s_i\in V$ 和目标 $g_i\in V$。时间离散为 $t=0,1,2,\dots$，每个时间步每台机器人执行一个动作：\textbf{移动}到相邻顶点，或\textbf{等待}（停在原地）。机器人 $i$ 的路径 $\pi_i=(\pi_i(0),\pi_i(1),\dots)$ 是它在各时刻所在的顶点序列，满足 $\pi_i(0)=s_i$、最终停在 $g_i$、相邻时刻要么相同（等待）要么是 $\mathcal{G}$ 中的一条边（移动）。一个 MAPF \textbf{解}是一组路径 $\{\pi_1,\dots,\pi_k\}$，两两之间既无顶点冲突也无边冲突。
\end{definition}

\subsection{两类冲突：顶点冲突与边冲突}
\label{sec:mapf-conflicts}

\begin{definition}[顶点冲突与边冲突]\label{def:mapf-conflict-types}
两条路径之间的两类冲突\cite{stern2019mapf}：
\begin{itemize}
  \item \textbf{顶点冲突}（vertex conflict）$\langle i,j,v,t\rangle$：两台机器人在同一时刻 $t$ 占据同一顶点 $v$，即 $\pi_i(t)=\pi_j(t)=v$；
  \item \textbf{边冲突 / 交换冲突}（edge / swap conflict）$\langle i,j,u,v,t\rangle$：两台机器人在 $t\to t{+}1$ 之间交换位置，即 $\pi_i(t)=\pi_j(t{+}1)=u$ 且 $\pi_i(t{+}1)=\pi_j(t)=v$，对应“两车在一条边上对向相遇穿模”。
\end{itemize}
\end{definition}

漏掉边冲突是初学者最常见的错误——只查“同一格”会让对向交换的两台机器人“擦肩穿过”。除了顶点和边，Stern 等人 $2019$ 的基准还梳理了几种更严格的碰撞模型\cite{stern2019mapf}：\textbf{跟随冲突}（following，机器人 $i$ 在 $t{+}1$ 进入 $j$ 在 $t$ 刚离开的格子）、\textbf{循环冲突}（cycle，三台及以上首尾相接同步轮转）、\textbf{膨胀冲突}（把机器人当成占多个格的圆盘）。

\begin{table}[H]\centering\small
\caption{MAPF 冲突类型谱}
\label{tab:mapf-conflict-spectrum}
\begin{tabular}{llll}
\toprule
冲突类型 & 何时发生 & 物理含义 & 何时需要禁止 \\
\midrule
顶点冲突 & 两机同时刻占同一格 & 两机撞在一点 & 总是（最基本） \\
边/交换冲突 & 两机同时刻对穿一条边 & 两车对向穿模 & 总是（最基本） \\
跟随冲突 & $i$ 进入 $j$ 刚离开的格 & 紧贴车尾跟车 & 担心刹车不及/惯性大时 \\
循环冲突 & 三机及以上首尾相接同步轮转 & 一圈机器人集体转位 & 要求严格无环依赖时 \\
膨胀/几何冲突 & 两机间距小于安全距离 & 占多格的圆盘相刮 & 机器人有真实尺寸时 \\
\bottomrule
\end{tabular}
\end{table}

经典 MAPF 只取前两类（点机器人 $+$ 同步假设）；后三类对应真实尺寸、刹车距离与连续运动，要么以附加约束加进搜索，要么交给 \cref{sec:mapf-exec} 的连续时间几何模型处理。

\subsection{两种优化目标：sum-of-costs 与 makespan}
\label{sec:mapf-objectives}

\begin{definition}[SOC 与 makespan]\label{def:mapf-objectives}
无碰撞的解通常有很多，两种主流目标：
\begin{itemize}
  \item \textbf{代价和}（sum-of-costs, SOC）：$\sum_{i} \mathrm{cost}(\pi_i)$，其中 $\mathrm{cost}(\pi_i)$ 是机器人 $i$ 到达并停在目标所用的时间步数。优化总吞吐；
  \item \textbf{完工时间}（makespan）：$\max_i \mathrm{cost}(\pi_i)$，即最后一台机器人到达目标的时刻。优化“全部完成”的时间。
\end{itemize}
\end{definition}

二者一般给出不同的最优解，且\textbf{都是 NP-难}\cite{yu2013structure}。CBS 默认优化 SOC。举个让二者分道扬镳的小例子：三台机器人因共享一条窄道，无法都走最短路。方案甲——让一台绕远（代价 $10$）、另两台走最短（各 $2$）：$\mathrm{SOC}=14$、$\mathrm{makespan}=10$。方案乙——让三台均摊绕行（各 $6$）：$\mathrm{SOC}=18$、$\mathrm{makespan}=6$。SOC 偏好甲（$14<18$，总吞吐更优），makespan 偏好乙（$6<10$，最后一台更早完成）。直觉：SOC 鼓励“牺牲少数、成全多数”，makespan 鼓励“负载均衡、不让任何一台拖后腿”。

\subsection{耦合 vs 解耦：一条连续谱}
\label{sec:mapf-spectrum}

把 MAPF 算法摆在“耦合$\leftrightarrow$解耦”的谱上看最清楚：

\begin{itemize}
  \item \textbf{完全耦合}（联合 A\*）：在联合状态空间搜索，分支因子 $\sim b^k$（$b$ 为单体分支因子）。最优、完备，但指数爆炸，$k>10$ 基本不可行；
  \item \textbf{完全解耦}（独立规划 $+$ 优先级）：各机器人单独 A\*，靠优先级/预留表避让（\cref{sec:mapf-scalable}）。快、可扩展，但一般\textbf{不完备、不最优}；
  \item \textbf{CBS（\cref{sec:mapf-cbs}）居中}：搜索是单体的（低层单体 A\*），但用“约束树”逐步消解冲突来恢复最优性——兼得解耦的搜索效率与耦合的最优保证。
\end{itemize}

理解了这条谱，后面所有 MAPF 算法你都能一眼定位：它在哪个位置、拿什么换什么。

\begin{center}
\small
\begin{tikzpicture}[>=Stealth]
  \draw[->, thick, gray!60!black] (0,0) -- (11,0);
  \node[anchor=north, font=\footnotesize] at (0.6,-0.15) {完全解耦};
  \node[anchor=north, font=\footnotesize] at (10.4,-0.15) {完全耦合};
  \node[draw, rounded corners, fill=second!10, font=\footnotesize, align=center] (pp) at (1.6,0.9) {独立 A\* / PP\\ 快·不完备};
  \node[draw, rounded corners, fill=main!10, font=\footnotesize, align=center] (cbs) at (5.5,0.9) {CBS\\ 单体搜索·最优};
  \node[draw, rounded corners, fill=third!10, font=\footnotesize, align=center] (joint) at (9.4,0.9) {联合 A\*\\ 最优·$b^k$ 爆};
  \draw[->, gray] (pp.south) -- (1.6,0.05);
  \draw[->, gray] (cbs.south) -- (5.5,0.05);
  \draw[->, gray] (joint.south) -- (9.4,0.05);
\end{tikzpicture}
\end{center}

\subsection{代码验证：冲突检测与联合搜索的爆炸}
\label{sec:mapf-def-code}

先把两类冲突的检测写清楚（这是 \cref{sec:mapf-cbs} CBS 的基础工具），再用联合 A\* 在极小实例上亲眼看分支爆炸。

\begin{codebox}[Python]{冲突检测与联合 A\* 的分支爆炸}
import numpy as np
from itertools import product
from heapq import heappush, heappop

def pad(path, T):                                   # 把路径补齐到长度 T(到目标后原地等待)
    return list(path) + [path[-1]] * (T - len(path))

def find_first_conflict(paths):
    """检测第一个冲突: ('vertex',i,j,v,t) 或 ('edge',i,j,u,v,t) 或 None。"""
    T = max(len(p) for p in paths); P = [pad(p, T) for p in paths]; k = len(P)
    for t in range(T):
        for i in range(k):                          # 顶点冲突
            for j in range(i + 1, k):
                if P[i][t] == P[j][t]: return ('vertex', i, j, P[i][t], t)
        if t + 1 < T:                               # 边冲突(交换)
            for i in range(k):
                for j in range(i + 1, k):
                    if P[i][t] == P[j][t + 1] and P[i][t + 1] == P[j][t]:
                        return ('edge', i, j, P[i][t], P[i][t + 1], t)
    return None

def neighbors(grid, p):
    r, c = p; R, C = grid.shape; out = [(r, c)]     # 等待
    for dr, dc in [(-1, 0), (1, 0), (0, -1), (0, 1)]:
        nr, nc = r + dr, c + dc
        if 0 <= nr < R and 0 <= nc < C and grid[nr, nc] == 0: out.append((nr, nc))
    return out

def joint_astar(grid, starts, goals, max_expand=2_000_000):
    """完全耦合的联合 A*:状态=所有机器人位置的元组。演示分支爆炸。"""
    k = len(starts)
    def h(s): return sum(abs(s[i][0]-goals[i][0])+abs(s[i][1]-goals[i][1]) for i in range(k))
    start = tuple(starts); goal = tuple(goals)
    openh = [(h(start), 0, start)]; gbest = {start: 0}; expanded = 0
    while openh:
        f, g, s = heappop(openh); expanded += 1
        if expanded > max_expand: return None, expanded
        if s == goal: return g, expanded
        per_agent = [neighbors(grid, s[i]) for i in range(k)]
        for moves in product(*per_agent):           # 联合后继:笛卡尔积
            if len(set(moves)) != k: continue       # 顶点冲突
            if any(moves[i] == s[j] and moves[j] == s[i]
                   for i in range(k) for j in range(i + 1, k)): continue   # 交换冲突
            ns = tuple(moves); ng = g + k
            if ng < gbest.get(ns, 1e9):
                gbest[ns] = ng; heappush(openh, (ng + h(ns), ng, ns))
    return None, expanded

grid = np.zeros((4, 4), dtype=int)
corners = [(0, 0), (3, 3), (0, 3), (3, 0), (1, 1)]
for k in range(2, 6):
    starts = corners[:k]; goals = [corners[(i + k // 2) % k] for i in range(k)]
    cost, expanded = joint_astar(grid, starts, goals)
    print(f"k={k}: 联合 A* 展开 {expanded:>9,} 个状态 | 解代价 {cost}")
\end{codebox}

\textbf{运行要点}：展开的状态数随机器人数 $k$ \textbf{超指数式}增长——这就是“把多机当一个联合智能体搜”的代价，联合分支因子 $\sim 5^k$。\verb@find_first_conflict@ 则是后面 CBS 的核心工具：它必须同时查顶点冲突\textbf{和}边冲突（注意 $t{+}1$ 那段），漏掉边冲突会让对向交换的机器人“穿模”而检测不出。把这两段代码吃透，你就握有了 MAPF 的“问题语言”。

\begin{pitfall}[概念误区：只检测顶点冲突，漏掉边/交换冲突]
判断两条路径无碰撞时，只检查“是否在同一时刻占同一格”，不检查“是否在一条边上对向交换”。后果：两台对向机器人在一条边上交换位置被判为“无冲突”，实际执行时正面对撞或穿模。根本原因：碰撞有两类（\cref{def:mapf-conflict-types}），只查前者漏掉后者。正确做法：冲突检测必须同时查两类（见代码 \verb@find_first_conflict@ 的 $t{+}1$ 分支）。
\end{pitfall}

\begin{pitfall}[思维陷阱：把 MAPF 当成“$N$ 个独立单体 A\*”]
认为多机路径规划就是每台机器人各跑一次 A\*、拼起来就行。后果：得到的路径集大概率含冲突（尤其拥挤场景），执行即死锁。根本原因：MAPF 的难度全部来自“无碰撞”这个机器人间的耦合约束；去掉耦合它就退化成 $N$ 个平凡的单体最短路。耦合正是它 NP-难的根源（\cref{def:mapf-objectives}）。正确做法：从一开始就把它当作有耦合约束的联合问题对待——要么 CBS、要么优先级规划/PIBT，而不是 $N$ 个独立 A\* 拼凑。
\end{pitfall}

\begin{pitfall}[编程陷阱：SOC 与 makespan 混用，或到达目标后等待的代价算错]
把 SOC（各路径长度之和）和 makespan（最长路径长度）混为一谈；或机器人到目标后继续等待时把等待也计入代价。后果：不同目标给出不同最优解，混用导致“最优解”对不上；到达目标后的等待若计入，会错误惩罚那些早到却需为他人让路的机器人。根本原因：$\mathrm{cost}(\pi_i)$ 的标准定义是“到达并\textbf{留在}目标所需的时间步”，到达后停在目标通常不再计费。正确做法：明确选定目标（CBS 默认 SOC）并在代码里统一；计算单体代价时只算到“最后一次离开目标后再不离开”的时刻。
\end{pitfall}

\begin{exercise}
\begin{enumerate}
  \item \textbf{[实现]} 扩展 \verb@find_first_conflict@ 让它返回\textbf{所有}冲突的列表，统计随机路径里顶点冲突与边冲突各占多少，并观察两类冲突比例随拥挤程度的变化。
  \item \textbf{[实现/分析]} 用 \verb@joint_astar@ 在固定 $4\times 4$ 地图上把机器人数从 $2$ 加到 $6$，记录展开状态数与运行时间，在对数坐标上拟合其增长阶，估算到 $k=?$ 就跑不动了。
  \item \textbf{[推导]} 构造一个 $2$ 机器人小例子，其 SOC 最优解的 makespan 不是最小、且 makespan 最优解的 SOC 不是最小（理解两个目标的内在冲突）。
  \item \textbf{[实现]} 实现一个“解的合法性验证器”：检查每条路径起于 $s_i$、终于 $g_i$，相邻时刻是等待或合法边，无顶点冲突，无边冲突（\cref{sec:mapf-cbs,sec:mapf-scalable} 都要反复用它）。
  \item \textbf{[思考，接 \cref{sec:mapf-cbs}]} CBS 的思路是先让每台机器人独立规划，发现冲突后只针对冲突“打补丁”。预测：这种策略在什么情况下会很快（冲突少时）、什么情况下仍可能退化（冲突极多、约束树爆炸）？
\end{enumerate}
\end{exercise}

% ════════════════════════════════════════════════════════════════
\section{CBS：冲突基搜索}
\label{sec:mapf-cbs}

\paragraph{动机：要最优，又不想爆炸。}
我们想要 MAPF 的\textbf{最优}解（像联合 A\* 那样），但联合 A\* 的 $5^N$ 分支让它在 $N>10$ 就死。关键观察是：\textbf{大多数时候，大多数机器人的路径其实互不干涉}——真正冲突的往往只是少数几对机器人在少数几个位置。既然如此，何必把所有机器人的决策耦合进一个巨大的联合空间？不如让每台机器人先独立规划，只在\textbf{真正发生冲突}的地方，才施加最小的额外约束去消解它。CBS 正是这个思路的精确实现。

\paragraph{反面：朴素“独立规划 $+$ 打补丁”会丢最优。}
发现机器人 $i,j$ 在顶点 $v$、时刻 $t$ 冲突，于是“让其中一个等一步”。但\textbf{让哪一个等？} 如果只武断地让 $i$ 等，可能错过“让 $j$ 绕路反而总代价更小”的更优解——这就丢了最优性。这正是优先级规划（\cref{sec:mapf-scalable}）的做法：固定一个优先级，低优先级的给高优先级的让路；它快，但\textbf{不完备也不最优}。CBS 的关键改进是：面对一个冲突，它\textbf{不武断选边}，而是\textbf{两种消解都试}——开两个分支，一个约束 $i$ 不能在 $(v,t)$、另一个约束 $j$ 不能在 $(v,t)$，分别重规划、分别算代价，在高层按代价做最优优先搜索。

\paragraph{历史：Sharon 等人的 CBS 及其变体谱系。}
CBS 由 Sharon、Stern、Felner、Sturtevant 于 $2015$ 年发表在 \emph{Artificial Intelligence}\cite{sharon2015cbs}，迅速成为最优 MAPF 的标准算法。此后衍生出一整个谱系：\textbf{ICBS}\cite{boyarski2015icbs} 优先消解“关键冲突”（cardinal conflict，即消解它必然使两个子节点代价都上升的冲突）以剪枝；\textbf{CBSH/CG/WDG}\cite{felner2018adding,li2019cbsh} 给高层搜索加入可采纳启发，大幅加速；\textbf{ECBS}\cite{barer2014ecbs} 放宽到有界次优。这些变体的共同点是都保留 CBS 的两层骨架，只在“选哪个冲突分裂”“高层如何估代价”上做文章。

\subsection{两层结构}
\label{sec:mapf-cbs-twolayer}

\begin{definition}[CBS 的高层与低层]\label{def:mapf-cbs-layers}
\begin{itemize}
  \item \textbf{高层：约束树（Constraint Tree, CT）}。CT 是一棵二叉树，每个节点 $N$ 含三样东西：一组\textbf{约束} $N.\mathrm{constraints}$（形如“机器人 $i$ 不能在时刻 $t$ 占顶点 $v$”或“……不能在 $t$ 走边 $(u,v)$”）、一组满足这些约束的\textbf{路径} $N.\mathrm{solution}$（每台机器人一条）、以及代价 $N.\mathrm{cost}=\mathrm{SOC}(N.\mathrm{solution})$。高层在 CT 上做\textbf{最优优先搜索}（按 cost 出队）；
  \item \textbf{低层：带约束的单体 A\*}。给定某机器人的约束集，低层为它单独搜一条满足所有约束的最短路（时空 A\*：状态是“(顶点, 时刻)”，避开被约束的时空点）。低层是\textbf{单体}的——这是 CBS 不爆炸的根本。
\end{itemize}
\end{definition}

\subsection{算法流程}
\label{sec:mapf-cbs-algo}

\begin{algo}[CBS（sum-of-costs）]\label{algo:mapf-cbs}
\begin{enumerate}
  \item \textbf{根节点}：无约束，每台机器人各跑一次普通 A\*（忽略他人），得到初始路径（可能含冲突）；
  \item \textbf{取出代价最小的 CT 节点 $N$}（最优优先），用 \verb@find_first_conflict@ 检验 $N.\mathrm{solution}$：
    \begin{itemize}
      \item \textbf{无冲突} $\Rightarrow$ $N$ 是目标节点，返回 $N.\mathrm{solution}$，\textbf{这就是最优解}；
      \item \textbf{有冲突} $\langle i,j,v,t\rangle$ $\Rightarrow$ \textbf{分裂}：生成两个子节点 $N_i,N_j$。$N_i$ 在 $N$ 的约束上\textbf{追加} $\langle i,v,t\rangle$，$N_j$ 追加 $\langle j,v,t\rangle$。对每个子节点，\textbf{只重规划被新约束的那台机器人}（低层 A\*），其余路径不变，更新代价，入高层队列；
    \end{itemize}
  \item 重复第 $2$ 步，直到取出无冲突节点。
\end{enumerate}
边冲突 $\langle i,j,u,v,t\rangle$ 的分裂类似：一个子节点禁 $i$ 走边 $(u,v)$ at $t$，另一个禁 $j$ 走 $(v,u)$ at $t$。
\end{algo}

\subsection{为什么 CBS 最优且完备}
\label{sec:mapf-cbs-optimal}

\begin{theorem}[CBS 的最优性与完备性]\label{thm:mapf-cbs}
CBS 返回的首个无冲突节点是 sum-of-costs 全局最优解，且 CBS 完备\cite{sharon2015cbs}。三点合起来保证最优：
\begin{enumerate}
  \item \textbf{低层最优}——每台机器人在给定约束下走的是满足约束的最短路；
  \item \textbf{高层最优优先 $+$ 代价可采纳}——CT 节点的代价 SOC 是“在该节点约束下任何无冲突解的代价下界”，最优优先搜索因此像 A\* 一样首次弹出的目标节点即全局最优；
  \item \textbf{分裂穷尽两种消解}——对每个冲突，“$i$ 让”和“$j$ 让”两种可能都被保留为子节点，不会因武断选边而漏掉最优解。
\end{enumerate}
完备性同理：若存在解，CT 的某个叶子必对应它。代价是最坏情况下 CT 节点数随冲突数指数增长。
\end{theorem}

\subsection{与优先级规划的对比}
\label{sec:mapf-cbs-vs-pp}

\begin{table}[H]\centering\small
\caption{CBS 与优先级规划（PP）对比}
\label{tab:mapf-cbs-pp}
\begin{tabular}{lll}
\toprule
维度 & CBS & 优先级规划（PP） \\
\midrule
最优性 & 最优（SOC） & 一般次优 \\
完备性 & 完备 & 不完备（优先级不当可能无解） \\
速度/可扩展 & 中（冲突少时快，多时爆） & 快，可扩展到大规模 \\
思路 & 冲突两种消解都试 & 固定优先级，低让高 \\
\bottomrule
\end{tabular}
\end{table}

一句话：CBS 用“两个子节点都展开”买来了最优与完备，代价是约束树可能爆炸；PP 用“武断选边”换来了速度，代价是丢掉最优与完备。\cref{sec:mapf-scalable} 会看到 PIBT/LaCAM 如何在 PP 这一侧把可扩展性推到极致。

\subsection{手把手走查一棵约束树}
\label{sec:mapf-cbs-trace}

用最经典的“对穿”实例：$3\times3$ 空栅格，机器人 $0$ 从 $(1,0)$ 去 $(1,2)$，机器人 $1$ 从 $(1,2)$ 去 $(1,0)$——它们必须穿过同一行。

\begin{itemize}
  \item \textbf{根节点}：无约束，各自跑最短路。$(1,0)\to(1,2)$ 唯一的两步最短路是 $(1,0)\to(1,1)\to(1,2)$，$(1,2)\to(1,0)$ 同理走 $(1,1)$。两条路都在 $t{=}1$ 占用中点 $(1,1)$ $\Rightarrow$ 顶点冲突 $\langle 0,1,(1,1),1\rangle$。根节点代价 $\mathrm{SOC}=2+2=4$；
  \item \textbf{分裂成两个子节点}，各加一条约束：
    \begin{itemize}
      \item \textbf{子节点 A}：约束“机器人 $0$ 在 $t{=}1$ 不得在 $(1,1)$”。只重规划机器人 $0$，它绕行 $(1,0)\to(0,0)\to(0,1)\to(0,2)\to(1,2)$（代价 $4$），机器人 $1$ 不变（代价 $2$），$\mathrm{SOC}=6$；
      \item \textbf{子节点 B}：约束“机器人 $1$ 在 $t{=}1$ 不得在 $(1,1)$”。对称地让机器人 $1$ 绕行，$\mathrm{SOC}=6$；
    \end{itemize}
  \item \textbf{高层按代价最优优先}：根（$\mathrm{SOC}=4$）弹出后压入两个 $\mathrm{SOC}=6$ 的子节点；弹出其中一个，检测到无冲突即返回——最优解 $\mathrm{SOC}=6$。
\end{itemize}

注意三个关键点：一，\textbf{只重规划被约束的那台机器人}（另一台路径原样保留），这是 CBS 比联合搜索高效的核心；二，\textbf{两个子节点都进队列}，不武断决定“谁让谁”，这是最优性的来源；三，\textbf{代价单调不减}——加约束只会让路径变长或不变，所以根的 $\mathrm{SOC}$ 是全局下界，最优优先才成立。本章代码在这个实例上报告展开 $4$ 个 CT 节点（而非最简故事里的 $2$ 个）——多出的来自“等待”产生的等代价或更低代价、但引入了\textbf{交换冲突}的中间路径，CBS 会继续对这些冲突分裂，直到弹出真正无冲突的叶子。

\subsection{CBS 的主要变体：都在跟“约束树爆炸”作斗争}
\label{sec:mapf-cbs-variants}

\begin{table}[H]\centering\small
\caption{CBS 主要变体（共同点：保留两层骨架，只把约束树修小）}
\label{tab:mapf-cbs-variants}
\begin{tabular}{llll}
\toprule
变体 & 加了什么 & 收益 & 仍最优？ \\
\midrule
ICBS\cite{boyarski2015icbs} & 优先消解 cardinal 冲突（分裂后两支代价都必涨） & 更快逼近无冲突 & 是 \\
CBSH/CG/DG/WDG\cite{felner2018adding,li2019cbsh} & 给高层 A\* 加可采纳启发（冲突图/依赖图） & 高层展开节点大减 & 是 \\
对称性消解（mutex 等） & 识别并一次性排除对称等价冲突 & 避免对称子树重复劳动 & 是 \\
绕行（bypassing） & 冲突可等代价改路绕开时不分裂、直接换路 & 省下大量分裂节点 & 是 \\
ECBS/EECBS（\cref{sec:mapf-bounded}） & 高低层改焦点搜索，放弃最优换有界次优 & 规模与速度大增 & 否（有界次优） \\
\bottomrule
\end{tabular}
\end{table}

前五种保持最优、只把树修小，第六种（\cref{sec:mapf-bounded} 详述）则干脆放松最优性换规模。CBS 框架的强大，正在于这些改进都能\textbf{正交地叠加}进“约束树 $+$ 低层搜索”这套骨架而不破坏其正确性。

\subsection{代码验证：CBS 求最优无冲突解}
\label{sec:mapf-cbs-code}

下面实现一个最小但完整的 CBS：低层是带时空约束的 A\*，高层在约束树上做最优优先搜索。复用 \cref{sec:mapf-def-code} 的 \verb@find_first_conflict@ 与 \verb@neighbors@。

\begin{codebox}[Python]{最小完整 CBS：低层时空 A\* $+$ 高层约束树}
from heapq import heappush, heappop
import itertools

def astar_st(grid, start, goal, vc, ec, max_t=64):
    """时空 A*: vc=顶点约束集{(v,t)}, ec=边约束集{(u,v,t)}(t->t+1 禁走 u->v)。
    返回到达并可停留在 goal 的最短路径(顶点序列)。"""
    last_goal_con = max([t for (v, t) in vc if v == goal], default=-1)   # 目标上最晚的约束
    def h(p): return abs(p[0] - goal[0]) + abs(p[1] - goal[1])
    openh = [(h(start), 0, start, 0)]; came, g_best = {}, {(start, 0): 0}
    while openh:
        f, g, p, t = heappop(openh)
        if p == goal and t > last_goal_con:                # 到目标且其后无约束逼它离开
            path = [p]
            while (p, t) in came: p, t = came[(p, t)]; path.append(p)
            return path[::-1]
        for np_ in neighbors(grid, p):
            nt = t + 1
            if (np_, nt) in vc: continue                   # 顶点约束
            if (p, np_, t) in ec: continue                 # 边约束(t 时刻从 p 走到 np_)
            ns = (np_, nt)
            if nt <= max_t and g + 1 < g_best.get(ns, 1e9):
                g_best[ns] = g + 1; came[ns] = (p, t)
                heappush(openh, (g + 1 + h(np_), g + 1, np_, nt))
    return None

def soc(paths): return sum(len(p) - 1 for p in paths)

def cbs(grid, starts, goals, max_nodes=100000):
    """最优 CBS(sum-of-costs)。返回(无冲突路径集, 代价, 展开的CT节点数)。"""
    k = len(starts)
    def low_level(constraints):
        paths = []
        for i in range(k):
            vc = {(v, t) for c in constraints if c[0] == 'v' for (_, a, v, t) in [c] if a == i}
            ec = {(u, v, t) for c in constraints if c[0] == 'e' for (_, a, u, v, t) in [c] if a == i}
            p = astar_st(grid, starts[i], goals[i], vc, ec)
            if p is None: return None
            paths.append(p)
        return paths
    root_paths = low_level(set()); counter = itertools.count()
    openh = [(soc(root_paths), next(counter), set(), root_paths)]; expanded = 0
    while openh:
        cost, _, constraints, paths = heappop(openh); expanded += 1
        if expanded > max_nodes: return None, None, expanded
        conf = find_first_conflict(paths)
        if conf is None: return paths, cost, expanded       # 首个无冲突节点即最优
        if conf[0] == 'vertex':
            _, i, j, v, t = conf; children = [('v', i, v, t), ('v', j, v, t)]
        else:
            _, i, j, u, v, t = conf; children = [('e', i, u, v, t), ('e', j, v, u, t)]
        for con in children:
            new_con = set(constraints); new_con.add(con); new_paths = low_level(new_con)
            if new_paths is not None:
                heappush(openh, (soc(new_paths), next(counter), new_con, new_paths))
    return None, None, expanded

grid = np.zeros((3, 3), dtype=int)
starts = [(1, 0), (1, 2)]; goals = [(1, 2), (1, 0)]         # 两机互换
paths, cost, expanded = cbs(grid, starts, goals)
print(f"CBS: SOC={cost} | 展开CT节点 {expanded} | 无冲突={find_first_conflict(paths) is None}")
\end{codebox}

\textbf{运行要点}：两台机器人要在 $3\times 3$ 地图上互换位置，直行会在中心格交换冲突；CBS 通过约束树发现冲突、分裂出“机器人 $0$ 不走中心边 / 机器人 $1$ 不走中心边”两个子节点，最终让其中一台绕行一格，返回\textbf{无冲突且 SOC 最优}的解。你可以把它和 \cref{sec:mapf-def-code} 联合 A\* 在同一实例上对比：CBS 的搜索全程是单体 A\*（低层），却拿到了和联合 A\* 一样的最优解——这就是“耦合的最优性 $+$ 解耦的效率”。

\begin{pitfall}[概念误区：冲突只展开一个子节点（只让一方避让）]
发现冲突后只生成“让机器人 $i$ 避让”这一个子节点，不生成“让 $j$ 避让”的另一个。后果：退化成了优先级规划——丢掉最优性（可能“让 $j$”才是更优解），甚至丢掉完备性（唯一可行解恰在被砍掉的那支）。根本原因：CBS 最优性的核心论证是“对每个冲突两种消解都保留为子节点”（\cref{thm:mapf-cbs}）；只展开一支等于武断选边。正确做法：每个冲突\textbf{必须}分裂出两个子节点（顶点冲突约束 $i$/约束 $j$；边冲突约束两个方向）。要加速用 ICBS 选 cardinal conflict、用 CBSH 加可采纳启发，而不是砍子节点。
\end{pitfall}

\begin{pitfall}[思维陷阱：低层 A\* 只避当前冲突点，不遵守该节点的全部约束]
在子节点重规划时，只让被约束的机器人避开“这次新加的”那个时空点，忽略它从祖先节点继承下来的其它约束。后果：重规划出的路径违反祖先约束，导致已消解的旧冲突复活，约束树陷入死循环、永不收敛。根本原因：CT 节点的约束是\textbf{沿树路径累积}的——子节点继承父节点的全部约束再加一条。正确做法：低层 A\* 的输入是该 CT 节点的\textbf{全部}约束（代码里 \verb@new_con = set(constraints); new_con.add(con)@），时空 A\* 扩展时对每条约束都检查。
\end{pitfall}

\begin{pitfall}[编程陷阱：时空 A\* 的目标判定忽略“目标点上的未来约束”]
低层 A\* 一旦机器人到达目标顶点就立即返回，不管目标点在更晚的时刻是否还有约束。后果：机器人“提前到达”并停在目标，但目标点在后续某时刻有约束（它得给路过的他人让出目标格），返回的路径其实违反约束。根本原因：MAPF 里机器人到达目标后要停留，若目标格在更晚时刻被约束，机器人必须先离开再回来或延后到达。正确做法：目标判定要求“到达目标且其后不再有针对该机器人在目标格的约束”（代码里 \verb@t > last_goal_con@）。
\end{pitfall}

\begin{exercise}
\begin{enumerate}
  \item \textbf{[实现/验证]} 跑通上面的 CBS，用 \cref{sec:mapf-def-code} 练习 4 的“合法性验证器”确认无冲突且起讫正确。再把同一实例喂给联合 A\*，验证两者 SOC \textbf{相同}（都最优，坐实 \cref{thm:mapf-cbs}）。
  \item \textbf{[实现/分析]} 在 $8\times8$ 地图上随机布置 $k\in\{2,\dots,8\}$ 台机器人，对每个 $k$ 跑 CBS 多次，记录展开 CT 节点数与求解时间，和联合 A\* 的曲线叠在一起——直观看到 CBS 增长比联合 A\* 平缓得多（但冲突极多时仍会上扬）。
  \item \textbf{[实现]} 给 CBS 加一个小优化：选择“分裂哪个冲突”时优先选 \textbf{cardinal conflict}（分裂后两子节点代价都严格上升）。对比加与不加时展开的 CT 节点数（这就是 ICBS 的核心思想）。
  \item \textbf{[实现/对比]} 实现一个优先级规划（给机器人固定优先级，按序用时空 A\* 规划、把已规划者的路径作为动态障碍），与 CBS 比较：(a) PP 快多少；(b) SOC 比 CBS 差多少；(c) 是否存在 CBS 有解而某些优先级顺序下 PP 无解的实例。
  \item \textbf{[推导]} 论证 \cref{thm:mapf-cbs}：为什么 CT 节点的代价是“该节点约束下任何无冲突解代价的下界”？为什么“高层最优优先 $+$ 两子节点都展开”就能保证首个弹出的无冲突节点是全局最优？
\end{enumerate}
\end{exercise}

% ════════════════════════════════════════════════════════════════
\section{可扩展 MAPF：优先级规划、PIBT 与 LaCAM}
\label{sec:mapf-scalable}

\paragraph{动机：当机器人有上千台、任务永不停止。}
CBS 最优、完备，但 \cref{sec:mapf-cbs} 末尾你已看到——冲突一多约束树就爆，它撑不到上千台机器人。可仓库里有上万台 AGV、空中有成百架无人机，而且任务\textbf{源源不断}（终身运行），根本等不起最优求解。CBS 被两个现实需求逼到墙角：其一，\textbf{规模}——亚马逊仓库里上万台机器人同时跑，CBS 的约束树会爆炸到无法求解；其二，\textbf{终身性}（lifelong）——仓库机器人送完一个包裹立刻领下一个目标，系统永远在跑。工程上需要的是\textbf{亚秒级出解、随规模近线性、能在线持续重规划}的方法，哪怕解不是最优——只要“足够好且无碰撞”。这就要回到 \cref{sec:mapf-cbs} 一开始放弃的那一侧：解耦/优先级。但纯独立规划会死锁，所以关键是设计\textbf{轻量却有协调机制}的方法。

\paragraph{反面：纯优先级规划不完备、还会活锁。}
最朴素的可扩展尝试就是“各自独立 A\*”——结果死锁。稍进一步，\textbf{优先级规划}（给机器人排个优先级、低的给高的让路）能消除大部分死锁，也确实快，但有两个固有缺陷：一是\textbf{不完备}——优先级排得不好，低优先级机器人可能根本找不到路（被高优先级机器人“停在目标上”堵死）；二是\textbf{不最优}——固定优先级武断地决定了谁让谁。更麻烦的是死锁的“升级版”——\textbf{活锁}（livelock）：在终身、滑动窗口的设定里，如果协调规则设计不当，机器人可能反复避让、互相礼让到谁也走不动。PIBT 用“优先级继承 $+$ 回溯”解决避让的局部死锁，LaCAM 在 PIBT 之上套一层搜索来恢复完备性——这是这条线的两个关键跃迁。

\paragraph{历史：从优先级规划到 PIBT、LaCAM。}
优先级规划是最老的多机路径思路，Erdmann 与 Lozano-P\'erez 早在 $1987$ 年就提出\cite{erdmann1987multiple}。$2005$ 年 Silver 的 Cooperative A\*（及其滑动窗口版 WHCA\*）把“预留表 $+$ 时空 A\*”做成游戏 AI 里的标准多机寻路\cite{silver2005cooperative}。真正的可扩展突破来自 Okumura 等人：$2019$ 年提出 \textbf{PIBT}（Priority Inheritance with Backtracking, IJCAI 2019，期刊版 $2022$ 年发表于 \emph{Artificial Intelligence}）\cite{okumura2022pibt}，一个极轻量的“单步” MAPF 规划器，能在亚毫秒内为上千台机器人生成下一步的无碰撞配置。$2023$ 年 Okumura 在 AAAI 提出 \textbf{LaCAM}（lazy constraints addition search for MAPF）\cite{okumura2023lacam}，用惰性后继生成把搜索做到既完备又能扩展到上千智能体；同年 IJCAI 的 \textbf{LaCAM\textsuperscript{*}} 加入 anytime 机制使它最终收敛到最优\cite{okumura2023lacamstar}。一个标志性事件：$2023$ 年首届 MAPF 竞赛（League of Robot Runners）的 $825$ 份提交中，夺冠策略的核心正是 PIBT。

\subsection{优先级规划（PP）与预留表}
\label{sec:mapf-pp}

PP 的机制：给 $k$ 台机器人定一个优先级顺序，\textbf{按序}逐台规划。规划第 $i$ 台时，把所有已规划的高优先级机器人的路径当作\textbf{时空障碍}——它们占用的每个“(顶点, 时刻)”都被记进\textbf{预留表}（reservation table），第 $i$ 台用时空 A\* 在“未被预留的时空”里找最短路。

\begin{insight}[预留表必须记三样东西]
预留表必须同时记\textbf{顶点占用}（防顶点冲突）、\textbf{边占用/反向禁止}（防交换冲突），还要处理“机器人到达目标后停在那儿”的占用（否则后续机器人会撞上停着的它）。第三样正是 PP 不完备的来源：高优先级机器人若停在咽喉要道（目标恰在唯一通道上），低优先级机器人可能永远绕不过去。
\end{insight}

PP 的代价：\textbf{不完备}（如上）与\textbf{不最优}（固定优先级武断决定避让方式）。Silver 的 WHCA\* 用\textbf{滑动窗口}（只规划未来 $w$ 步、滚动重规划）来降低单次规划量、适应动态，但窗口太短会引入活锁\cite{silver2005cooperative}。PP 快、实现简单、在不太拥挤的场景里够用，是一切可扩展 MAPF 的起点。

\subsection{PIBT：优先级继承与回溯}
\label{sec:mapf-pibt}

PIBT 把视角缩到“单步”：给定当前所有机器人的配置（各自所在顶点），它生成\textbf{下一时刻}的一个无碰撞、可转移的配置——即一个“one-step MAPF”求解器，反复调用就得到完整路径\cite{okumura2022pibt}。它的三个关键机制：

\begin{itemize}
  \item \textbf{自适应优先级}：机器人的优先级随“等待时间”动态上升（越久没到目标、优先级越高），保证每台机器人最终都能获得高优先级、被照顾到；
  \item \textbf{优先级继承}（priority inheritance）：高优先级机器人先选它最想去的相邻顶点；若那个顶点被一台低优先级机器人占着，这台低优先级机器人就\textbf{继承}高优先级，被“推着”也得给前者让位、去找自己的下一步；
  \item \textbf{回溯}（backtracking）：若被推的机器人发现自己无处可去（所有相邻顶点都不可行），它向上回传一个“无效”信号，触发上层机器人改选次优顶点。
\end{itemize}

\begin{theorem}[PIBT 的完备性条件]\label{thm:mapf-pibt}
当环境图满足“\textbf{任意一对相邻顶点都属于某个简单环}”（例如双连通图）时，无论多少台机器人，所有机器人都能在有限时间内到达各自目标\cite{okumura2022pibt}。直觉：简单环提供了“绕行让路”的回旋余地，不会出现死胡同顶死。
\end{theorem}

PIBT 极轻量——每次配置生成对上千台机器人也只要亚毫秒，这正是它能撑起 $2023$ 竞赛冠军的原因。但它是\textbf{短视的}（只看一步），整体路径质量不保证最优。

\paragraph{走一遍优先级继承。}
取一条三格走廊 $A\!-\!B\!-\!C$。机器人 $1$（高优先级）当前在 $A$，目标在 $C$，这一步最想去 $B$；机器人 $2$（低优先级）当前正占着 $B$。PIBT 按优先级从高到低决策：(1) 机器人 $1$ 先挑，它最想去的 $B$ 被机器人 $2$ 占着；(2) 触发\textbf{优先级继承}——机器人 $2$ 临时获得机器人 $1$ 的高优先级，被要求“让位”，它必须为自己找一个不与已决策者冲突的下一步；(3) 机器人 $2$ 的相邻格是 $A$（机器人 $1$ 正打算离开它去 $B$）和 $C$；为避免与机器人 $1$ 的目标格 $B$ 冲突、也避免占住 $A$ 造成对穿，机器人 $2$ 选择移到 $C$；(4) 机器人 $2$ 让出 $B$ 后，机器人 $1$ 顺利移入 $B$。最终这一步：机器人 $1$ 走 $A\to B$、机器人 $2$ 走 $B\to C$，无碰撞，且高优先级的机器人 $1$ 朝目标推进了一格。再看\textbf{回溯}何时发生：若机器人 $2$ 的所有相邻格都不可行，它就向上回传“无效”信号；机器人 $1$ 收到后放弃 $B$、改选次优动作（如原地等待）。

\subsection{LaCAM：惰性约束添加搜索}
\label{sec:mapf-lacam}

PIBT 虽快却短视、且其完备性依赖图的环结构。LaCAM\cite{okumura2023lacam} 在 PIBT 之上套一层搜索，换来\textbf{完备性}：它是一个搜索式算法（结构上类似 A\*），完备、次优，核心创新是\textbf{惰性后继生成}（lazy successor generation）——本质是一个“在配置空间上的惰性深度优先搜索”。LaCAM 的每个搜索节点存三样东西：(i) 一个\textbf{配置}（所有机器人当前所在顶点的组合）、(ii) 一棵用队列表示的\textbf{约束树}（决定从该配置生成后继时优先尝试哪些“低层约束”）、(iii) 指向\textbf{父节点}的指针（用于回溯出最终路径）。生成后继配置时，LaCAM 调用 \textbf{PIBT 作为配置生成器}——PIBT 负责快速产出一个满足当前约束的、无碰撞的下一配置。“惰性”体现在：不一次性枚举一个配置的所有后继，而是按约束树一点点、按需地生成。LaCAM 的完备性来自：只要给足时间，它最终会搜遍整个状态空间——因此若有解必能找到。

\subsection{LaCAM\textsuperscript{*} 与终身 MAPF}
\label{sec:mapf-lacam-star}

\textbf{LaCAM\textsuperscript{*}}\cite{okumura2023lacamstar} 给 LaCAM 加上 \textbf{anytime} 机制：它先快速给出一个可行解，然后随着时间推移不断改进，\textbf{最终收敛到最优}（当解的代价定义为“累积转移代价”时）。在 MAPF 标准基准（$33$ 张四连通栅格、共 $13900$ 个实例、机器人数最多 $1000$）上，LaCAM\textsuperscript{*} 解出的实例数远超此前方法。由于 MAPF 最优求解是 NP-难（\cref{def:mapf-objectives}），“先可行、再 anytime 逼近最优”是大规模实例的务实路线。

\textbf{终身 MAPF}（lifelong MAPF, LMAPF）是另一个工程维度：机器人到达目标后立刻被分配\textbf{新目标}，系统持续运行（正是仓库取送货的常态）。它需要\textbf{在线、滚动地重规划}：常用做法是周期性地用快速求解器（如 PIBT/LaCAM）对当前活跃目标重新求解一小段时间窗。LMAPF 把 \cref{sec:mapf-cbba} 的任务分配（不断来的新任务派给谁）与本节的路径规划（派定后怎么走）\textbf{耦合}在一个永不停止的循环里——这也是 \cref{sec:mapf-submod} 要讨论的“分配与路径如何接起来”在终身设定下的极致形态。

\begin{table}[H]\centering\small
\caption{四种可扩展 MAPF 方法对比（从纯解耦逐步加回协调/搜索）}
\label{tab:mapf-scalable-compare}
\begin{tabular}{llllll}
\toprule
方法 & 核心机制 & 完备性 & 最优性 & 单步/整体开销 & 典型定位 \\
\midrule
优先级规划 PP & 固定优先级 $+$ 预留表，低让高 & 不完备 & 次优 & 一次规划，快 & 不拥挤场景的最简基线 \\
PIBT & 单步配置生成 $+$ 优先级继承 $+$ 回溯 & 双连通图上完备 & 短视次优 & 单步亚毫秒，上千台 & 实时、终身、超大规模底座 \\
LaCAM & 配置空间惰性 DFS，用 PIBT 作生成器 & 完备 & 次优 & 搜索，极省内存 & 要完备 $+$ 规模 \\
LaCAM\textsuperscript{*} & LaCAM $+$ anytime 改进 & 完备 & 最终最优 & anytime，随时可停 & 要完备 $+$ 规模 $+$ 渐进最优 \\
\bottomrule
\end{tabular}
\end{table}

这张表读下来有一条清晰主线：从 PP 到 LaCAM\textsuperscript{*}，是不断给“纯解耦”加回协调与搜索——PP 完全武断，PIBT 用继承/回溯加了局部协调（换来环上完备），LaCAM 在外面套一层搜索（换来全局完备），LaCAM\textsuperscript{*} 再加 anytime（换来渐进最优）。每加一层能力更强，但也更重——这又是本章那根“协调质量 vs 计算代价”权衡轴在可扩展这一侧的细分。

\subsection{代码验证：优先级规划 $+$ 预留表（看它的快与不完备）}
\label{sec:mapf-scalable-code}

下面实现 PP $+$ 预留表（复用 \cref{sec:mapf-cbs-code} 的 \verb@astar_st@、\verb@find_first_conflict@、\verb@soc@），在一个\textbf{瓶颈地图}上演示它的“顺序依赖 / 不完备”——同一实例，换个优先级顺序，结果从“无解”变“有解”，而 CBS 不依赖顺序、总能解。

\begin{codebox}[Python]{优先级规划 $+$ 预留表：快但不完备、对顺序敏感}
def prioritized_planning(grid, starts, goals, order=None, max_t=48):
    """PP+预留表。按 order 逐台规划,每台把已规划者的时空占用当障碍。
    返回 (路径集, None) 成功; 或 (None, 卡住的机器人) 失败——PP 不完备。"""
    k = len(starts)
    if order is None: order = list(range(k))
    res_v, res_e = set(), set()                          # 预留表:顶点(v,t)、边(v,u,t)禁反向
    paths = [None] * k
    for i in order:
        p = astar_st(grid, starts[i], goals[i], set(res_v), set(res_e), max_t)
        if p is None: return None, i                     # 第 i 台被堵死
        paths[i] = p
        for t, v in enumerate(p): res_v.add((v, t))
        for t in range(len(p), max_t + 1): res_v.add((p[-1], t))     # 到目标后永久占用
        for t in range(len(p) - 1):
            u, v = p[t], p[t + 1]; res_e.add((v, u, t))              # 阻止他人反向交换
    return paths, None

g = np.zeros((3, 3), dtype=int); g[1, 0] = 1; g[1, 2] = 1   # (1,1)是连接上下两行的唯一通道
starts = [(0, 0), (2, 0)]; goals = [(1, 1), (0, 2)]         # 机器人0的目标恰是瓶颈格
for order in [[0, 1], [1, 0]]:
    paths, stuck = prioritized_planning(g, starts, goals, order=order)
    if paths is None:
        print(f"PP order={order}: 失败! 机器人{stuck}被堵死(高优先级机器人停在瓶颈)")
    else:
        print(f"PP order={order}: 成功 SOC={soc(paths)} 无冲突={find_first_conflict(paths) is None}")
r = cbs(g, starts, goals, max_nodes=20000)                 # 同实例,CBS 不依赖顺序
print(f"CBS(同实例): SOC={r[1]} 展开{r[2]}个CT节点 (无论顺序都能解)")
\end{codebox}

\textbf{运行要点}：\verb@order=[0,1]@ 时，机器人 $0$（高优先级）规划到瓶颈格 $(1,1)$ 并永久停在那里，把唯一通道堵死，机器人 $1$ \textbf{找不到任何路径}——PP 失败。仅仅把顺序换成 \verb@order=[1,0]@，让机器人 $1$ 先过、机器人 $0$ 后到并适当等待，就\textbf{成功}了。而 CBS 在同一实例上不依赖任何顺序、总能找到无冲突最优解。这就把“PP 快但不完备、且对优先级顺序敏感”变成了你能复现的现象——也解释了为什么工程上用 PIBT（优先级继承让被堵的机器人能推开挡路者）、用 LaCAM（套一层搜索恢复完备性）来补 PP 的这两个洞。PIBT/LaCAM 的完整实现较长（可参考 Okumura 的 \verb@pypibt@ 最小 Python 实现），其机制已在上面理论部分讲清。

\begin{pitfall}[概念误区：以为优先级规划/PIBT 给的是最优解]
把 PP 或 PIBT 的输出当成最优 MAPF 解，用它替代 CBS 还期望同样的质量。后果：解的总代价明显比最优差（某些机器人绕了大远路），在质量敏感场景上吃亏却不自知。根本原因：PP/PIBT 都在“解耦/优先级”这一侧，用最优性换可扩展；PIBT 还是\textbf{短视}的（只规划一步）。正确做法：明确需求——要最优且规模不大用 CBS；要规模/实时可接受次优用 PIBT/LaCAM；要“先可行再逼近最优”用 LaCAM\textsuperscript{*}（anytime）。
\end{pitfall}

\begin{pitfall}[思维陷阱：预留表只记顶点占用，漏掉边（交换）与目标永久占用]
实现预留表时只记录“某顶点某时刻被占”，不记边的反向禁止，也不处理“机器人到目标后一直停在那儿”的占用。后果：后规划的机器人与已规划者发生交换冲突（穿模），或一头撞上停在目标格的高优先级机器人。根本原因：多机碰撞有顶点 $+$ 边两类（\cref{def:mapf-conflict-types}）；且机器人到达目标后通常停留，该格在后续所有时刻都被占用。正确做法：预留表三样都要记——顶点占用 \verb@(v,t)@、边反向禁止 \verb@(v,u,t)@、目标到达后对所有未来时刻的占用（代码里 \verb@for t in range(len(p), max_t+1)@）。
\end{pitfall}

\begin{pitfall}[编程陷阱：终身/滑动窗口里重规划不当导致活锁]
在终身 MAPF 里用很短的滑动窗口反复重规划，或每次重规划都重置优先级/不保留历史承诺。后果：机器人在路口反复“你让我、我让你”，谁也不前进——活锁。根本原因：窗口太短看不到死锁的“出路”，或优先级/承诺不稳定导致决策反复横跳。正确做法：用带完备性保证的求解器（LaCAM）做底，或保证优先级单调（PIBT 的自适应优先级让久等者优先级升高），滑动窗口要足够长以越过常见瓶颈，并对“长期不前进”的机器人做兜底处理（提权或全局重规划）。
\end{pitfall}

\begin{exercise}
\begin{enumerate}
  \item \textbf{[实现/分析]} 在上面 PP 代码基础上实现一个“\textbf{随机重启}”策略：PP 失败时随机打乱优先级顺序重试若干次。在多个含瓶颈的地图上统计“需要几次重启才成功”以及“成功解的 SOC 比 CBS 最优差多少”。
  \item \textbf{[实现]} 给 PP 加滑动窗口（只规划未来 $w$ 步、到窗口末尾再滚动重规划），做成一个简易 WHCA\*。在终身场景里观察窗口长度 $w$ 对吞吐与活锁的影响。
  \item \textbf{[实现/对比]} 实现一个 PIBT 的单步配置生成器（按自适应优先级 $+$ 优先级继承 $+$ 回溯，生成下一无碰撞配置），反复调用得到完整路径。在双连通栅格上跑上百台机器人，记录每步耗时（应为亚毫秒级）与到达率，和 CBS 在同规模上的求解时间对比。
  \item \textbf{[思考/设计]} PIBT 的完备性要求“任意相邻顶点对都在某个简单环上”（\cref{thm:mapf-pibt}）。举出一个\textbf{不满足}该条件的地图（如带死胡同的树状走廊），说明 PIBT 在上面可能让某些机器人永远到不了目标，再设计一个最小的地图改动（加一条边）使其变双连通。
  \item \textbf{[思考，接 \cref{sec:mapf-submod}]} 终身 MAPF 把“任务分配”和“路径规划”耦合进一个永不停止的循环。设计这个循环的架构：多久重分配一次任务、多久重规划一次路径、两者如何交接、当某机器人的路径长期走不通时如何反馈去重新分配。
\end{enumerate}
\end{exercise}

% ════════════════════════════════════════════════════════════════
\section{边界、子模视角与编队的对照}
\label{sec:mapf-submod}

\paragraph{动机：把碎片拼成一张完整的图。}
读到这里，你手里有一堆算法：匈牙利、拍卖、CBBA、CBS、PIBT、LaCAM。但真实系统不是单个算法——它要先分配（谁去做哪个任务）、再规划路径（怎么走不撞），而且这两步互相影响：分配时若不考虑路径会不会堵，可能派出一个“看着近、实际堵死”的方案；路径规划时分配已定，只能在既定目标下腾挪。同时你可能困惑：\cref{ch:mr-consensus} 讲的编队（连续、平滑）和本章的 MAPF（离散、栅格）看着像两个世界，它们到底什么关系？这一节把这些线索拧成一股。

\subsection{子模视角：$50\%$ 与 $1-1/e$ 背后的同一结构}
\label{sec:mapf-submod-theory}

回看 \cref{sec:mapf-sga,sec:mapf-cbba}：SGA/CBBA 在 DMG 条件下有 $50\%$ 界。这个“边际收益递减”其实就是组合优化里的\textbf{子模性}（submodularity）：

\begin{definition}[子模性]\label{def:mapf-submodular}
集合函数 $f$ 子模 $\iff$ 往小集合里加一个元素的增益不小于往它的超集里加同一元素的增益，即对 $S\subseteq T$ 且 $e\notin T$ 有 $f(S\cup\{e\})-f(S)\ge f(T\cup\{e\})-f(T)$。这正是 \cref{def:mapf-dmg} 的 DMG。
\end{definition}

子模性是多机“分工”问题能被贪婪高效逼近的根本原因，它有两个经典保证：

\begin{theorem}[子模贪婪的两个近似界]\label{thm:mapf-submod-bounds}
\begin{itemize}
  \item \textbf{基数约束}（最多选 $K$ 个）下，贪婪最大化单调子模函数，解 $\ge (1-1/e)\approx 0.63$ 倍最优（Nemhauser--Wolsey--Fisher 1978）\cite{nemhauser1978analysis}。覆盖、传感、信息增益这类多机任务效用大多是单调子模的；
  \item \textbf{拟阵约束}（matroid，如“每台机器人容量上限”这类更一般的结构）下，贪婪的保证降为 $1/2$——这正是 \cref{sec:mapf-cbba} CBBA 的 $50\%$ 界的来历（任务分配的约束是拟阵）。
\end{itemize}
\end{theorem}

把任务分配放进子模优化的框架，$1-1/e$ 与 $1/2$ 这些界就不再是孤立结论，而是同一棵树上的果子。\textbf{分布式子模}则把这套搬到多机：Corah 与 Michael（\emph{Autonomous Robots} 2019）用机器人间的顺序贪婪做分布式子模最大化并保持近似比\cite{corah2019distributed}；Xu、Garimella 与 Tzoumas（\emph{IEEE T-RO} 2025）进一步处理分布式与\textbf{抗失效}（部分机器人掉线仍有保证）的子模任务分配\cite{xu2025distributed}。\pz{存疑（suspect，cite\_key=xu2025distributed）：书目元数据准确（T-RO 41:3480--3499, 2025），但“抗失效/fault-tolerant”系误标——该文核心是 RAG（Resource-Aware distributed Greedy），主攻\emph{通信与计算高效}（决策时间随网络规模线性增长 vs SOTA 立方增长），全文不涉及 resilience/机器人失效。“抗失效”实对应同作者群另一篇\emph{不同}论文（Zhou--Tzoumas--Pappas--Tokekar, Distributed Attack-Robust Submodular Maximization, arXiv:1910.01208，非本引用）。建议：把本引用的性质改述为“通信/计算高效的分布式子模”；如需“抗失效”另引前述论文。}这条线是 \cref{ch:mr-consensus} “分布式优化”的\textbf{离散表亲}——共识章的 ADMM 切分连续优化，这里的分布式贪婪切分组合优化，内核都是“用局部计算 $+$ 有限协调逼近全局最优”。

\subsection{把“分配”与“路径”接起来}
\label{sec:mapf-submod-pipeline}

一个完整的多机调度系统通常是三段：\textbf{任务分配}（task $\to$ robot）$\to$ \textbf{指派/排序}（每台机器人的任务序列）$\to$ \textbf{路径规划}（MAPF，把序列变成无冲突时空路径）。难点在于这三段\textbf{互相耦合}：

\begin{itemize}
  \item 分配时\textbf{不知道}真实路径代价——你用直线距离估代价派活，但实际路径可能因拥堵远比直线长，导致“看着最优的分配，执行起来很堵”；
  \item 路径规划时分配\textbf{已经定死}——MAPF 只能在既定的起讫点下找无冲突路径，若分配本身就让两台机器人必然抢一条窄道，MAPF 再优也回天乏术。
\end{itemize}

务实的做法有几种：\textbf{串行}（先分配再规划，简单但易堵）、\textbf{迭代}（规划发现某分配走不通就反馈重分配）、\textbf{联合}（把分配与路径放进一个优化/搜索里同时解，最优但难）。终身 MAPF（\cref{sec:mapf-lacam-star}）更是把这个循环做成永不停止的在线过程。

\subsection{离散 MAPF 与连续编队的边界}
\label{sec:mapf-submod-boundary}

本章的 MAPF 和 \cref{ch:mr-consensus} 的编队，看似两个世界，实则是\textbf{同一个问题在不同抽象层的两种建模}——都在回答“多个机器人如何在共享空间里不相撞地完成各自目标”：

\begin{table}[H]\centering\small
\caption{离散 MAPF 与连续编队/控制的对照}
\label{tab:mapf-vs-formation}
\begin{tabular}{lll}
\toprule
维度 & 离散 MAPF（本章） & 连续编队/控制（\cref{ch:mr-consensus}$+$后续） \\
\midrule
空间/时间 & 图/栅格 $+$ 离散步 & 连续空间 $+$ 连续时间 \\
决策对象 & 走哪格、何时走、谁先过 & 控制量（速度/加速度/力） \\
处理什么 & 组合避让、路口排序 & 动力学、平滑、物理交互 \\
擅长场景 & 密集结构化（仓库巷道、路网） & 平滑机动、编队、协同搬运 \\
不擅长 & 动力学、力、连续协调 & 大规模离散排序、密集路口调度 \\
典型衔接 & 输出离散路点序列 & 跟踪这些路点（MPC 参考轨迹） \\
\bottomrule
\end{tabular}
\end{table}

二者常常\textbf{混合}：用 MAPF 在离散层规划出无冲突的\textbf{路点序列}（谁先过路口、谁走哪条巷），再用连续控制/MPC 去\textbf{平滑跟踪}这些路点（处理动力学、避免急停急转）。这里最该记住的一点：它们不是竞争关系，而是\textbf{抽象层次}关系——离散 MAPF 在“决策层”决定宏观走法，连续控制在“执行层”实现微观运动。

\subsection{代码验证：子模贪婪的 $1-1/e$ 界，与违反子模时的失效}
\label{sec:mapf-submod-code}

下面实现一个\textbf{覆盖类}任务分配：$K$ 台传感机器人，每台覆盖一组目标点，效用 $=$ 被覆盖的\textbf{不同}目标数（单调子模），贪婪每轮选“新增覆盖最多”的那台。在小规模上枚举最优，验证贪婪解恒 $\ge (1-1/e)$ 倍最优；再把效用换成“目标需被 $\ge 2$ 台\textbf{同时}覆盖才算数”（协同/超可加，\textbf{违反}子模），看保证如何崩掉。

\begin{codebox}[Python]{子模贪婪的 $1-1/e$ 界，与违反子模时的失效}
import numpy as np, math
from itertools import combinations

def make_instance(n_cand, n_targets, seed, p=0.25):
    rng = np.random.default_rng(seed)
    return [set(np.where(rng.random(n_targets) < p)[0]) for _ in range(n_cand)]

def coverage(selected, cover):                      # 单调子模效用:被覆盖的不同目标数
    s = set()
    for i in selected: s |= cover[i]
    return len(s)

def greedy_cover(cover, K):
    n = len(cover); sel = []
    for _ in range(K):                              # 每轮选边际增益最大的候选
        best, gain = None, -1
        for i in range(n):
            if i in sel: continue
            g = coverage(sel + [i], cover) - coverage(sel, cover)
            if g > gain: gain, best = g, i
        sel.append(best)
    return coverage(sel, cover)

def optimal_cover(cover, K):                         # 小规模枚举最优作对照
    return max(coverage(c, cover) for c in combinations(range(len(cover)), K))

bound = 1 - 1 / math.e                               # ≈ 0.632
ratios = [greedy_cover(c := make_instance(10, 20, s), 3) / optimal_cover(c, 3) for s in range(80)]
print(f"子模覆盖: 贪婪/最优 均值={np.mean(ratios):.3f} 最差={min(ratios):.3f} | 1-1/e={bound:.3f}")
print(f"  所有实例都 >= 1-1/e ? {min(ratios) >= bound - 1e-9}")
\end{codebox}

\textbf{运行要点}：子模覆盖下，贪婪解在 $80$ 个随机实例上的最差比值约 $0.87$、均值约 $0.99$，\textbf{全部 $\ge 1-1/e\approx0.63$}——这正是 \cref{thm:mapf-submod-bounds} 的实测体现（它是最坏界，典型实例远好于此）。而一旦把效用换成“需两台同时覆盖才算”的协同形式（边际收益变成\textbf{递增}而非递减，违反子模），贪婪的最差比值跌到 $0.5$、\textbf{击穿了 $1-1/e$ 界}。守住子模，贪婪就有常数比保证；违反子模（强协同任务），保证就失效——这也是 MR 协同搬运为何更难的根本原因。

\subsection{多视角、对比与本质洞察}
\label{sec:mapf-submod-views}

\textbf{三个看待本章的视角}：

\begin{itemize}
  \item \textbf{优化视角}：任务分配是匹配/整数规划（指派问题、子模最大化），MAPF 是带耦合约束的最短路问题。难度都来自“耦合”；
  \item \textbf{市场/经济视角}：拍卖、价格、竞标。任务是商品，机器人是买家，价格是协调信号——CBBA 把这个信号用共识在图上传播；
  \item \textbf{搜索视角}：CBS、LaCAM 本质都是树搜索——在“约束/配置”空间里做最优优先或惰性 DFS。
\end{itemize}

\begin{insight}[本质洞察一：减耦合换可扩展、保耦合换最优]
匈牙利/CBS 保住全局耦合 $\to$ 最优但成本爆；贪婪/优先级/PIBT 砍掉耦合 $\to$ 可扩展但次优。这和共识章“全局协调 vs 局部自主”、ADMM 的 $\rho$ 旋钮是\textbf{同一个权衡的不同化身}——多机系统反复在“协调的质量”与“计算/通信的代价”之间做这道选择题。认清这一点，本章十来个算法就不再是孤立技巧，而是同一根轴上的不同落点。
\end{insight}

\begin{insight}[本质洞察二：子模性是“分工”可被高效逼近的统一根源]
边际收益递减（子模性）让任务分配的 $50\%$ 界（拟阵）、覆盖/传感的 $1-1/e$ 界（基数）、CBBA 与集中贪婪的等价、分布式贪婪的近似比，统统从同一个性质里长出来。一旦效用违反子模（出现协同增益/超可加），这些保证就集体失效——这也是为什么“多机协作搬一件重物”（强协同、超可加）比“多机分头巡逻”（子模）在算法上难得多。
\end{insight}

\begin{pitfall}[概念误区：把子模性当成凸性或线性]
看到“贪婪有保证”就以为目标函数是凸的/线性的，套用连续优化的直觉。后果：误用连续优化方法或误判问题难度。根本原因：子模是\textbf{集合函数}上的“边际收益递减”，与连续函数的凸性是不同概念（虽有类比）；子模最大化是 NP-难但贪婪有常数比，凸优化是多项式可解。正确做法：先判断效用是不是单调子模（\cref{def:mapf-submodular}）；是则贪婪有 $1-1/e$（基数）或 $1/2$（拟阵）保证；不是（如协同增益）换专门方法。
\end{pitfall}

\begin{pitfall}[思维陷阱：分配与路径割裂优化，分配只看直线距离]
任务分配阶段用“机器人到任务的直线/曼哈顿距离”当代价求最优分配，完全不考虑执行时的拥堵与路径冲突。后果：分配在纸面上最优，一执行就发现多台机器人挤向同一窄道、实际完工时间远超预期。根本原因：真实代价是无冲突路径的长度，不是直线距离；在拥挤环境里两者差距巨大。正确做法：在拥挤场景里分配阶段就用更接近真实的代价估计（考虑拥堵的路径代价、或历史路径长度），或采用迭代/联合方法。
\end{pitfall}

\begin{pitfall}[编程陷阱：把任务分配的“通信图”和路径规划的“环境图”搞混]
实现时只维护一张图，把“机器人间谁能通信”（CBBA 用）和“环境里哪些位置相邻可走”（MAPF 用）混为一张。后果：CBBA 在错误的图上传播中标信息（收敛行为错乱），或 MAPF 在通信拓扑上找路（完全无意义）。根本原因：本章有\textbf{两张本质不同的图}——通信图（节点$=$机器人）与环境图/栅格图（节点$=$位置）。正确做法：代码里用不同的数据结构和命名清晰区分（如 \verb@comm_graph@ vs \verb@env_grid@）；CBBA 的收敛分析用通信图的直径，MAPF 的搜索用环境图。
\end{pitfall}

\begin{exercise}
\begin{enumerate}
  \item \textbf{[实现/推导]} 实现一个覆盖类任务的贪婪选择，验证贪婪解 $\ge (1-1/e)$ 倍最优（小规模枚举对比）。再构造一个违反子模的协同效用，观察贪婪保证失效。把 $1-1/e$ 界与 \cref{sec:mapf-cbba} 的 $50\%$ 界放一起，理解它们都源于子模。
  \item \textbf{[实现/分析，综合]} 搭一个最小的“分配 $+$ 路径”流水线：用匈牙利按\textbf{直线距离}分配任务，再用 CBS 规划无冲突路径，记录“分配阶段估计的总距离”与“CBS 实际总路径长度”的差距，再用“实际路径长度”迭代重分配一次看差距是否缩小。
  \item \textbf{[思考/设计]} 给定场景判断该用离散 MAPF、连续编队、还是混合方案并说明理由：(a) $50$ 台 AGV 在网格化仓库取送货；(b) $4$ 架无人机保持菱形编队巡线；(c) $6$ 台机器人先各自到达指定货位（密集巷道）、再两两协同把大件抬到打包区。
  \item \textbf{[推导]} 说明为什么“基数约束 $+$ 单调子模”贪婪是 $1-1/e$，而“拟阵约束”降为 $1/2$（\cref{thm:mapf-submod-bounds}）。提示：$1-1/e$ 来自每步贪婪至少拿到“剩余最优增益的 $1/K$”，拟阵情形用交换论证。
  \item \textbf{[思考，接后续]} 设想后续的分布式 MPC 编队要跟踪本章 MAPF 给出的离散路点序列，会遇到哪些离散到连续的衔接问题：路点间如何插值成平滑参考轨迹？机器人动力学约束与 MAPF 的“每步走一格”假设如何调和？当连续跟踪偏离离散计划时是否需要触发 MAPF 重规划？
\end{enumerate}
\end{exercise}

% ════════════════════════════════════════════════════════════════
\section{有界次优与 anytime 大规模 MAPF：ECBS/EECBS 与 MAPF-LNS}
\label{sec:mapf-bounded}

\paragraph{动机：要规模，也要对质量心里有数。}
把 \cref{sec:mapf-cbs,sec:mapf-scalable} 的方法摆出来看：CBS 给最优，但几十台、冲突一密就跑不动；PIBT 亚毫秒能上千台，但路径可能绕很多冤枉路、且不给任何“离最优多远”的保证；LaCAM\textsuperscript{*} 虽 anytime 最终最优，但“最终”可能很久。工程上反复出现的真实诉求是：“给我一个解，质量别太离谱（比如保证不超过最优的 $1.5$ 倍），而且要快、要能扩展。”这就是\textbf{有界次优}（bounded-suboptimal）的用武之地——用一个用户指定的次优因子 $w\ge 1$ 换速度，同时保证 $\mathrm{SOC}\le w\cdot\mathrm{SOC}^*$。另一条诉求是：“先给我个能用的，然后有多少时间就优化多少。”这就是 \textbf{anytime} 的思路。

\paragraph{反面：无界放松 vs 反复重启都不好。}
最朴素的“次优”做法：把 CBS 的最优优先随便放松一下——挑个“看着快收敛”的节点先展开。问题：这样得到的解\textbf{没有任何质量保证}，可能离最优很远，你却不知道有多远。次优不可怕，可怕的是\textbf{无界}的次优。所以关键不是“放松”，而是\textbf{有原则地放松}：放松到什么程度、对应多大的质量损失，必须可量化、可保证。另一种朴素的 anytime 做法：反复用不同随机种子从头跑 PP，留最好的（random restart）——它能改进，但每次都从零开始、丢掉之前的搜索成果，效率低。MAPF-LNS 的关键改进是：不重启，而是\textbf{在已有解上做局部手术}——只重规划一小撮智能体，保留其余，逐步把解改好。

\paragraph{历史：焦点搜索、ECBS、EECBS 与 LNS。}
有界次优的工具是\textbf{焦点搜索}（focal search, Pearl \& Kim 1982）\cite{pearl1982focal}：在最优优先的基础上，用一个次优因子框出一批“足够好”的节点，在其中按另一个启发优先展开。把它用到 MAPF，Barer 等人 $2014$ 年提出 \textbf{ECBS}（Enhanced CBS），在 CBS 的高层和低层都用焦点搜索，首个有界次优 MAPF 算法\cite{barer2014ecbs}。$2021$ 年 Li、Ruml 与 Koenig 提出 \textbf{EECBS}（Explicit Estimation CBS, AAAI 2021）\cite{li2021eecbs}，把高层的焦点搜索换成显式估计搜索（EES）并用\textbf{在线学习}估计每个 CT 节点解的代价，成为当时最快的有界次优求解器。anytime 这条线，Li 等人 $2021$ 年提出 \textbf{MAPF-LNS}（基于大邻域搜索的 anytime MAPF）\cite{li2021lns}，$2022$ 年的 \textbf{LNS2}（AAAI 2022）进一步能从\textbf{带碰撞}的初始解出发，通过修复把碰撞消到零\cite{li2022lns2}。

\subsection{焦点搜索与有界次优}
\label{sec:mapf-focal}

回忆 A\*/最优优先：用一个优先队列 OPEN 按 $f=g+h$ 排序，每次弹出 $f$ 最小的。\textbf{焦点搜索}在此之上加一个 FOCAL 列表：

\begin{equation}
  \mathrm{FOCAL}=\{\,n\in\mathrm{OPEN}\ :\ f(n)\le w\cdot f_{\min}\,\},\qquad f_{\min}=\min_{n\in\mathrm{OPEN}} f(n).
  \label{eq:mapf-focal}
\end{equation}

即所有 $f$ 值不超过当前最小 $f$ 的 $w$ 倍的节点。算法\textbf{从 FOCAL 里}按另一个（可不可采纳的）启发 $h_2$（比如“这个节点的解里还有多少冲突”）挑一个展开。

\begin{theorem}[焦点搜索的有界次优性]\label{thm:mapf-focal}
由于 $f_{\min}$ 始终是最优解代价的下界，而我们只展开 $f\le w\cdot f_{\min}$ 的节点，\textbf{最终解的代价保证 $\le w\cdot\mathrm{SOC}^*$}\cite{pearl1982focal,barer2014ecbs}。$w=1$ 时 FOCAL 退回 OPEN，就是最优搜索；$w$ 越大，FOCAL 越宽，越能用 $h_2$ 引导快速冲向目标，牺牲的质量上界也越大。
\end{theorem}

这就是“有原则地放松最优性”。

\subsection{ECBS 与 EECBS}
\label{sec:mapf-eecbs}

\textbf{ECBS} 把焦点搜索用在 CBS 的\textbf{两层}：高层 CT 上，FOCAL 里按“解中冲突数最少”优先展开（冲突少的节点更可能接近无冲突解）；低层单体 A\* 上，也用焦点搜索，在“足够短”的路径里挑\textbf{与他人冲突最少}的——这能直接减小 CT 的规模\cite{barer2014ecbs}。整体保证 $\mathrm{SOC}\le w\cdot\mathrm{SOC}^*$。

\textbf{EECBS} 指出 ECBS 的高层焦点搜索在两个启发（代价下界 vs 冲突数）负相关时会变低效，于是把高层换成\textbf{显式估计搜索}（EES）：它额外维护一个对“每个 CT 节点最终解代价”的\textbf{不可采纳估计}（用在线学习从已展开节点拟合），用这个更准的估计来决定先展开谁，同时仍用可采纳的代价下界保证 $w$-界\cite{li2021eecbs}。EECBS 配合绕行、冲突优先级、对称性消解等改进，显著快于 ECBS，是有界次优 MAPF 的代表。它的骨架仍是 \cref{sec:mapf-cbs} 的 CBS——约束树 $+$ 低层单体搜索——只是“选哪个节点展开”用了更聪明、且仍带保证的策略。

\subsection{MAPF-LNS：anytime 大邻域搜索}
\label{sec:mapf-lns}

MAPF-LNS 走完全不同的路子——\textbf{先有一个解，再改进}\cite{li2021lns}：

\begin{algo}[MAPF-LNS]\label{algo:mapf-lns}
\begin{enumerate}
  \item \textbf{初始解}：用一个快速方法（如优先级规划 PP，或 PIBT）拿一个可行但可能次优的解；
  \item \textbf{破坏}（destroy）：选一小撮智能体（一个“邻域”，如随机若干个、或冲突/绕路最多的若干个），把它们的路径删掉；
  \item \textbf{修复}（repair）：把这撮智能体的路径\textbf{重新规划}——把其余智能体的当前路径当作固定的时空障碍（预留表），用 PP 或低层搜索为这撮重规划；
  \item 若新解总代价更低，\textbf{接受}；否则保留旧解。回到第 $2$ 步，直到时间用完。
\end{enumerate}
\end{algo}

这是个 anytime 过程：任何时候停下都有一个合法解，且解随时间单调变好、逼近最优。邻域怎么选很关键（随机、基于冲突、基于代价、甚至用学习来选），直接决定改进速度。\textbf{LNS2} 把这套用于\textbf{修复可行性}：从一个带碰撞的初始解出发，反复选“还有碰撞的”智能体子集重规划以减少碰撞数，直到为零——这在超大规模、连可行解都难找的实例上特别有效\cite{li2022lns2}。

\subsection{把所有 MAPF 方法放进一张二维谱}
\label{sec:mapf-2dspectrum}

到这里，本章的 MAPF 方法可以摆进一张二维图——横轴“规模/速度”，纵轴“质量保证”：

\begin{table}[H]\centering\small
\caption{MAPF 方法的“质量保证 $\times$ 规模”二维谱}
\label{tab:mapf-2d}
\begin{tabular}{llll}
\toprule
方法 & 质量保证 & 规模/速度 & 代表 \\
\midrule
最优 & 最优（$\mathrm{SOC}^*$） & 小（冲突少时尚可） & CBS、CCBS \\
有界次优 & $\le w\cdot\mathrm{SOC}^*$ & 中到大 & ECBS、EECBS \\
anytime & 先可行、随时间 $\to$ 近最优 & 大 & MAPF-LNS、LaCAM\textsuperscript{*} \\
无界次优/规则 & 无保证（但通常够用） & 极大、实时 & PP、PIBT、LaCAM、LNS2（修复） \\
\bottomrule
\end{tabular}
\end{table}

选型就是在这张图上按“我能接受多少次优 $\times$ 我要多大规模/多快”落点。这也呼应了贯穿本章的那根权衡：质量与代价的取舍，只是这里把“质量”精确成了“离最优的可证明倍数”。

\subsection{代码验证一：ECBS-lite——有界次优用 $w$ 换速度}
\label{sec:mapf-ecbs-code}

先实现一个 ECBS-lite：在 \cref{sec:mapf-cbs-code} CBS 的高层换成焦点搜索（FOCAL $=$ 代价 $\le w\cdot$ 当前最小代价的 CT 节点，其中选\textbf{冲突最少}者展开），保证 $\mathrm{SOC}\le w\cdot$ 最优。复用前文的 \verb@astar_st@、\verb@soc@、\verb@find_first_conflict@、\verb@pad@。

\begin{codebox}[Python]{ECBS-lite：高层焦点搜索，有界次优}
import heapq, itertools

def count_conflicts(paths):
    """一组路径的冲突总数(顶点+边)——焦点搜索的副启发(越少越接近无冲突)。"""
    T = max(len(p) for p in paths); P = [pad(p, T) for p in paths]; k = len(P); n = 0
    for t in range(T):
        for i in range(k):
            for j in range(i + 1, k):
                if P[i][t] == P[j][t]: n += 1
                if t + 1 < T and P[i][t] == P[j][t + 1] and P[i][t + 1] == P[j][t]: n += 1
    return n

def ecbs_lite(grid, starts, goals, w=1.5, max_nodes=100000):
    """有界次优 CBS:高层焦点搜索。OPEN 按代价排,FOCAL={cost<=w*最小代价}内选冲突最少者。"""
    k = len(starts)
    def low_level(cons):
        ps = []
        for i in range(k):
            vc = {(v, t) for c in cons if c[0] == 'v' for (_, a, v, t) in [c] if a == i}
            ec = {(u, v, t) for c in cons if c[0] == 'e' for (_, a, u, v, t) in [c] if a == i}
            p = astar_st(grid, starts[i], goals[i], vc, ec)
            if p is None: return None
            ps.append(p)
        return ps
    root = low_level(set()); cnt = itertools.count()
    OPEN = [(soc(root), next(cnt), frozenset(), root)]; heapq.heapify(OPEN); expanded = 0
    while OPEN:
        fmin = OPEN[0][0]                                   # 当前最小代价 = 最优下界
        focal = [nd for nd in OPEN if nd[0] <= w * fmin]    # FOCAL: 代价在 w 倍界内
        node = min(focal, key=lambda nd: count_conflicts(nd[3]))   # 选冲突最少者
        OPEN.remove(node); heapq.heapify(OPEN)
        cost, _, cons, paths = node; expanded += 1
        if expanded > max_nodes: return None, None, expanded
        conf = find_first_conflict(paths)
        if conf is None: return paths, cost, expanded       # 首个无冲突节点,SOC<=w*最优
        if conf[0] == 'vertex':
            _, i, j, v, t = conf; ch = [('v', i, v, t), ('v', j, v, t)]
        else:
            _, i, j, u, v, t = conf; ch = [('e', i, u, v, t), ('e', j, v, u, t)]
        for c in ch:
            nc = set(cons); nc.add(c); np_ = low_level(nc)
            if np_ is not None:
                heapq.heappush(OPEN, (soc(np_), next(cnt), frozenset(nc), np_))
    return None, None, expanded

g = np.zeros((10, 10), dtype=int)                          # 10x10,中间一道带缺口的竖墙
for r in range(2, 8):
    if r not in (4, 5): g[r, 4] = 1
s  = [(2, 7), (5, 4), (4, 4), (0, 3), (7, 3), (5, 9), (7, 7), (8, 7)]   # 8 台机器人
go = [(7, 1), (0, 6), (9, 6), (2, 2), (4, 8), (0, 1), (1, 5), (6, 1)]
opt = cbs(g, s, go, max_nodes=30000)
print(f"CBS 最优: SOC={opt[1]} 展开 {opt[2]} 个 CT 节点")
for w in [1.0, 1.3, 1.8]:
    p, c, e = ecbs_lite(g, s, go, w=w)
    print(f"  ECBS-lite w={w}: SOC={c} 展开{e:>5} | SOC<=w*最优={c <= w * opt[1] + 1e-9}")
\end{codebox}

\textbf{运行要点}：这个实例里 CBS 求最优要展开\textbf{两千多}个 CT 节点（冲突密、约束树爆）。而 ECBS-lite 只要放松一点点——$w{=}1.3$——就只展开\textbf{十几个}节点拿到只比最优多 $1$ 的解，提速约\textbf{两个数量级}，且解始终满足 $\mathrm{SOC}\le w\cdot$ 最优（\cref{thm:mapf-focal}）。这把本节的核心思想变成了你能在屏幕上读到的数字：\textbf{用可证明的微小质量损失，换取巨大的速度提升}——这正是大规模 MAPF 里有界次优远比死磕最优实用的原因。

\subsection{代码验证二：MAPF-LNS——anytime 改进}
\label{sec:mapf-lns-code}

下面在一张带短墙（制造拥堵）的 $6\times6$ 栅格上，先用优先级规划拿初始解，再用 MAPF-LNS 破坏-修复地改进，看 SOC 如何 anytime 地降到 CBS 最优。复用前文的 \verb@astar_st@、\verb@find_first_conflict@、\verb@soc@、\verb@prioritized_planning@、\verb@cbs@。

\begin{codebox}[Python]{MAPF-LNS：破坏-修复 anytime 改进}
def replan_subset(grid, starts, goals, paths, subset, max_t=50):
    """把 subset 之外的智能体当固定时空障碍,为 subset 重规划(LNS 的'修复')。"""
    res_v, res_e = set(), set()
    for i, p in enumerate(paths):
        if i in subset: continue
        for t, v in enumerate(p): res_v.add((v, t))
        for t in range(len(p), max_t + 1): res_v.add((p[-1], t))     # 目标永久占用
        for t in range(len(p) - 1):
            u, v = p[t], p[t + 1]; res_e.add((v, u, t))              # 禁反向交换
    new = {}
    for i in subset:                                                  # subset 内顺序规划
        p = astar_st(grid, starts[i], goals[i], set(res_v), set(res_e), max_t)
        if p is None: return None
        new[i] = p
        for t, v in enumerate(p): res_v.add((v, t))
        for t in range(len(p), max_t + 1): res_v.add((p[-1], t))
        for t in range(len(p) - 1):
            u, v = p[t], p[t + 1]; res_e.add((v, u, t))
    return new

def mapf_lns(grid, starts, goals, init_paths, iters=200, nbr=4, seed=0):
    """anytime 大邻域搜索:反复'破坏一小撮 + 修复',接受更优解。"""
    rng = np.random.default_rng(seed); k = len(starts)
    paths = [list(p) for p in init_paths]; best = soc(paths); hist = [best]
    for _ in range(iters):
        subset = list(rng.choice(k, size=min(nbr, k), replace=False)); rng.shuffle(subset)
        new = replan_subset(grid, starts, goals, paths, subset)
        if new is not None:
            cand = [new.get(i, paths[i]) for i in range(k)]
            if find_first_conflict(cand) is None and soc(cand) < best:   # 更优才接受
                paths = cand; best = soc(cand)
        hist.append(best)
    return paths, best, hist

g = np.zeros((6, 6), dtype=int); g[3, 1] = g[3, 2] = g[3, 3] = 1     # 带缺口的短墙
s  = [(1, 0), (4, 5), (0, 2), (0, 1), (4, 4), (5, 0)]; go = [(3, 0), (1, 2), (5, 2), (4, 0), (2, 3), (0, 0)]
pp, _ = prioritized_planning(g, s, go, order=list(range(6)), max_t=50)   # 初始解
fp, fs, hist = mapf_lns(g, s, go, pp, iters=200, nbr=4, seed=3)
opt = cbs(g, s, go, max_nodes=30000)[1]
print(f"PP 初始 SOC={soc(pp)} -> LNS 最终 SOC={fs} (CBS 最优={opt})")
print(f"SOC 改进轨迹(每25轮): {[hist[i] for i in range(0, 201, 25)]}")
\end{codebox}

\textbf{运行要点}：优先级规划给出的初始解 SOC 偏高（本例 $48$，因为拥堵地图下固定优先级让一些机器人绕了大圈），MAPF-LNS 在\textbf{数十轮}破坏-修复内就把它压到 $35$——恰好等于 CBS 的最优解。改进轨迹是单调下降的（anytime 的标志）：任何时刻停下都有合法解，给的时间越多越接近最优。这正是大规模 MAPF 的务实哲学——不追求一步到位的最优，而是“快速可行 $+$ 持续改进”。

\begin{pitfall}[概念误区：把“有界次优”和“无界次优/无保证”混为一谈]
看到 ECBS/EECBS“次优”就以为和 PIBT 一样没质量保证，或反过来以为随便放松 CBS 就有界。后果：向上游错误承诺质量，或放着有界次优不用、在不需要最优时硬扛最优 CBS 的爆炸。根本原因：有界次优有可证明的上界 $\mathrm{SOC}\le w\cdot\mathrm{SOC}^*$（来自焦点搜索的 $f_{\min}$ 下界，\cref{thm:mapf-focal}）；随意放松最优优先则没有任何界。正确做法：要质量保证就用焦点搜索类方法（ECBS/EECBS）并明确写出 $w$；只要“够用、极快”且能接受无保证才用 PIBT/LaCAM。
\end{pitfall}

\begin{pitfall}[思维陷阱：LNS 的邻域（破坏哪些智能体）随便选]
每轮都纯随机选一撮智能体重规划，不看哪些智能体真的“有改进空间”。后果：大量轮次花在重规划“本来就最优”的智能体上，改进极慢。根本原因：LNS 的改进速度高度依赖邻域质量——重规划那些\textbf{绕路多、与他人冲突多}的智能体，才最可能降代价。正确做法：用有信息的邻域选择——基于“当前路径比其单体最短路长多少”、基于冲突图的连通分量、或学习式选择。随机只作基线。
\end{pitfall}

\begin{pitfall}[编程陷阱：LNS 修复时没把其余智能体当完整的时空障碍]
重规划子集时只避开其余智能体当前所在的格子，忽略它们的边占用、以及到达目标后的永久占用。后果：修复出的新路径与“固定”的其余智能体产生新的边冲突或撞上停在目标的智能体，接受后整体解反而非法。根本原因：修复 $=$ 在“其余智能体路径固定”的子问题上重规划，必须包含顶点、边、目标永久占用三类预留。正确做法：\verb@replan_subset@ 里对每个非子集智能体登记完整时空占用，并在接受前用 \verb@find_first_conflict@ 复核整体无冲突。
\end{pitfall}

\begin{exercise}
\begin{enumerate}
  \item \textbf{[实现/分析]} 把 MAPF-LNS 的邻域选择从“随机”改成“优先重规划当前路径比其单体最短路长得最多的智能体”，在同一拥堵实例上对比两种邻域选择的 SOC 改进曲线，验证有信息的邻域收敛更快。
  \item \textbf{[实现]} 扫描 ECBS-lite 的 $w\in\{1.0,1.1,1.5,2.0\}$，记录展开 CT 节点数与解的 SOC，验证：$w$ 越大越快、SOC 越高，且始终 $\le w\cdot$ 最优（\cref{thm:mapf-focal}）。
  \item \textbf{[实现/对比]} 在 $w=1.5$ 下对比 ECBS-lite 与 MAPF-LNS 在一批中等规模实例上的“达到该质量所需时间”，体会有界次优（一次给出带保证的解）与 anytime（持续改进、随时可停）的不同适用场景。
  \item \textbf{[实现，接 LNS2]} 改造 MAPF-LNS 为“修复可行性”模式：从一个带碰撞的初始解出发，每轮选“仍有碰撞的”智能体子集重规划以减少碰撞总数，直到无碰撞（这就是 LNS2 的核心思想）。
  \item \textbf{[推导]} 证明 \cref{thm:mapf-focal}：为什么“只展开 $f\le w\cdot f_{\min}$ 的节点”能保证返回解的代价 $\le w\cdot\mathrm{SOC}^*$？提示：$f_{\min}$ 在搜索全程是最优解代价的下界。
\end{enumerate}
\end{exercise}

% ════════════════════════════════════════════════════════════════
\section{从规划到执行：连续时间、运动学约束与鲁棒执行}
\label{sec:mapf-exec}

\paragraph{动机：一个完美的离散计划，上了真机却撞了。}
前面所有方法都默认了三个理想化假设——时间离散成等长步、机器人是无尺寸的点、所有机器人\textbf{同步}执行。设想你用 CBS 求出了一个完美的离散 MAPF 计划——每个时间步谁在哪格、绝无冲突。可一旦放到真机上：机器人 A 因为轮子打滑慢了半拍，机器人 B 照原计划准时到达路口，结果两者撞在一起。离散 MAPF 的“无冲突”保证\textbf{建立在“所有机器人同步、每步等长”的假设上}；真机一旦有延迟、有速度差异，这个时间对齐就崩了。更何况真机不是点——它有尺寸、有运动学约束（差速轮机器人不能瞬间掉头），离散网格的“每步走一格”根本没刻画这些。

\paragraph{反面：加大时间余量 vs 全局重规划都不好。}
最朴素的鲁棒执行：给每步都留出大段时间余量、让机器人慢慢走。问题：\textbf{极度浪费}——为了应付偶发延迟，所有机器人全程龟速，吞吐暴跌；而且余量定多大全凭拍脑袋，延迟一旦超过余量照样撞。另一种朴素做法：延迟发生时整个系统\textbf{全局重规划}。问题：重规划开销大、频繁触发会让系统卡顿。真正的解法是\textbf{结构化}地处理执行不确定性：要么在规划时就把延迟容忍度\textbf{显式建模}进去（$k$-robust），要么把计划\textbf{后处理}成一个不依赖绝对时间、只依赖\textbf{相对顺序}的执行策略（动作依赖图）——只要保持“谁先过某个路口、谁后过”这个顺序，无论各机器人快慢、延迟多少，都不会撞。

\paragraph{历史：CCBS、MAPF-POST/ADG、k-robust 与 MAPD。}
连续时间这条线：Phillips 与 Likhachev $2011$ 年的 \textbf{SIPP}（安全区间路径规划）能在动态障碍中求时间最优路径\cite{phillips2011sipp}；Andreychuk 等人 $2019$ 年（期刊版 $2022$ 于 AIJ）提出 \textbf{CCBS}（连续时间 CBS），把 CBS 框架搬到连续时间 $+$ 几何体智能体\cite{andreychuk2022ccbs}。执行鲁棒性这条线：H\"onig 等人 $2016$ 年（ICAPS）提出 \textbf{MAPF-POST}\cite{honig2016mapfpost}，$2019$ 年（RA-L）的\textbf{动作依赖图}（Action Dependency Graph, ADG）进一步把计划编码成动作间的依赖顺序\cite{honig2019persistent}。Atzmon 等人 $2018$ 年（SoCS）提出 \textbf{$k$-robust MAPF}\cite{atzmon2018krobust}。终身这条线：Ma 等人 $2017$ 年（AAMAS）提出 \textbf{MAPD}（终身取送货）与 \textbf{Token Passing}\cite{ma2017lifelong}，Ma 与 Koenig $2016$ 年（AAMAS）提出 \textbf{TAPF / CBM}（联合目标分配与路径）\cite{ma2016optimal}。

\subsection{连续时间 MAPF 与 CCBS}
\label{sec:mapf-ccbs}

经典 MAPF 假设时间离散、每个动作恰好一步。\textbf{MAPF\textsubscript{R}}（连续时间 MAPF）放开这两点：时间连续、动作时长可不同（走对角线比走直线久）、智能体有几何形状。\textbf{CCBS} 是它的最优解法，沿用 CBS 的两层骨架但做两处关键改造\cite{andreychuk2022ccbs}：低层用 \textbf{SIPP}——它在“(顶点, 安全时间区间)”的状态空间里搜索，把动态障碍（其他智能体）留下的“不安全区间”挖掉，直接求连续时间下的时间最优路径；高层的冲突变成“两个带时刻的动作 $(a_i,t_i)$ 与 $(a_j,t_j)$ 在几何上相撞”，约束则指定“某动作必须避开的不安全时间区间”。CCBS 证明了在连续时间下仍保持最优与完备。

\begin{insight}[SIPP 如何把“离散步”跨到“连续时间”]
普通时空 A\* 的状态是 $(\text{顶点}, t)$，$t$ 取整数步——若时间连续，$t$ 有无穷多取值，状态空间炸成无穷。SIPP 的洞察：对一个顶点，真正重要的不是“哪个精确时刻能到”，而是“哪些\textbf{连续时间区间}里它是安全的（无动态障碍占用）”\cite{phillips2011sipp}。它把每个顶点的时间轴按其他智能体的占用切成若干\textbf{安全区间}（safe interval），状态就压缩成 $(\text{顶点}, \text{第几个安全区间})$——有限个！例如顶点 $v$ 被一个动态障碍在时间 $[3,5]$ 占用，则 $v$ 的安全区间是 $[0,3)$ 与 $(5,\infty)$ 两段；SIPP 搜索时机器人到达 $v$ 只需记录“落在哪个安全区间”，并在该区间内取\textbf{最早可行到达时刻}。这样连续时间被有限个区间“离散化”了，却没有丢掉任何时间精度。
\end{insight}

\subsection{执行鲁棒性：从时间戳到依赖顺序}
\label{sec:mapf-adg}

离散计划给的是“绝对时间戳”（第 $t$ 步谁在哪），这在真机延迟下很脆。两类工具让执行鲁棒：

\begin{itemize}
  \item \textbf{MAPF-POST}\cite{honig2016mapfpost}：把离散计划后处理成一个\textbf{时序调度}（plan-execution schedule）。它用一个简单时序网络（STN）编码动作间的时间约束，产出能在\textbf{非完整约束}机器人上执行、考虑运动学限制（最大平移/转向速度）、并保证机器人间\textbf{安全距离}的执行方案——多项式时间完成；
  \item \textbf{动作依赖图（ADG）}\cite{honig2019persistent}：更彻底地抛弃绝对时间，只保留\textbf{相对顺序}。ADG 把每个机器人的每个动作作为节点，加两类依赖边：\textbf{同机顺序}（机器人 $i$ 的第 $k$ 个动作必须在第 $k{+}1$ 个之前）与\textbf{跨机顺序}（若计划中机器人 $i$ 比 $j$ 先经过某格 $v$，则“$i$ 离开 $v$”必须先于“$j$ 进入 $v$”）；
  \item \textbf{$k$-robust MAPF}\cite{atzmon2018krobust}：换个思路，在\textbf{规划时}就预留延迟容忍——要求计划满足“即使任意智能体累计延迟不超过 $k$ 步，也不会碰撞”。$k$ 越大越鲁棒、但代价越高。它和 ADG 互补：前者规划时建模延迟，后者执行时吸收延迟。
\end{itemize}

\begin{theorem}[ADG 执行的鲁棒无碰撞性]\label{thm:mapf-adg}
执行时，只要每个机器人按\textbf{与 ADG 一致的任意顺序}推进（谁慢了就等它，绝不抢在依赖前面），则\textbf{无论各机器人速度如何、延迟多久，都保证无碰撞}\cite{honig2019persistent}。这把“安全”从脆弱的时间对齐，变成了鲁棒的顺序约束。
\end{theorem}

\subsection{终身取送货 MAPD 与联合目标分配-路径 TAPF}
\label{sec:mapf-mapd}

把 \cref{sec:mapf-scalable} 的终身 MAPF 落到仓库的具体形态，就是 \textbf{MAPD}（Multi-Agent Pickup and Delivery）：任务不断到达，每个任务是“去取货点取货、再送到送货点”；机器人完成一个立刻领下一个。\textbf{Token Passing}（TP）是其经典分布式解法——用一个“令牌”在机器人间传递，持有令牌的机器人为自己挑一个任务并规划路径（把其他机器人已规划的路径当约束），它能解所有“良构”（well-formed）实例\cite{ma2017lifelong}。MAPD 正是 \cref{sec:mapf-submod} 那个“任务分配 $\leftrightarrow$ 路径规划耦合”在\textbf{终身、在线}下的完整形态。

当任务/目标是\textbf{可互换的}（anonymous，谁去都行），问题变成 \textbf{TAPF}（联合目标分配与路径，Combined Target-Assignment and Path-Finding）。它的优雅解法 \textbf{CBM}（Conflict-Based Min-cost-flow）把 CBS 的高层约束树保留，但低层不再是单体 A\*——而是一个\textbf{最小费用最大流}：在“机器人-目标”二分结构上同时解出“谁去哪个目标”和“对应的无冲突路径”\cite{ma2016optimal}。CBM 是 \cref{sec:mapf-submod} “分配与路径如何接起来”这个问题最干净的\textbf{算法级}回答。

\begin{table}[H]\centering\small
\caption{从离散计划到真机：各工具的作用层与机制}
\label{tab:mapf-exec-tools}
\begin{tabular}{llll}
\toprule
工具 & 作用层 & 解决的问题 & 关键机制 \\
\midrule
CCBS & 规划层 & 连续时间、几何体、动作时长不一 & CBS $+$ SIPP 安全区间 $+$ 连续时间冲突 \\
MAPF-POST & 规划$\to$执行 & 把离散计划变成可执行运动学调度 & 简单时序网络，保证安全距离 \\
动作依赖图 ADG & 执行层 & 延迟/速度差下仍无碰撞 & 同机 $+$ 跨机相对顺序，谁慢等谁 \\
$k$-robust MAPF & 规划层 & 预留延迟容忍 & 时空缓冲，容忍至多 $k$ 步延迟 \\
MAPD / Token Passing & 终身调度 & 持续到达的取送货任务 & 令牌传递，解良构实例 \\
TAPF / CBM & 规划层 & 目标匿名时联合分配-路径 & CBS 高层 $+$ 最小费用流低层 \\
\bottomrule
\end{tabular}
\end{table}

一条清晰的分工：规划层（CCBS/$k$-robust/CBM）管“算出什么样的计划”，执行层（ADG/MAPF-POST）管“怎么把计划安全跑出来”，终身层（MAPD/TP）管“任务源源不断时怎么持续转”。

\subsection{代码验证：动作依赖图（ADG）让执行对延迟鲁棒}
\label{sec:mapf-adg-code}

下面用 CBS 解出的 $3\times 3$ 互换计划，构建 ADG，然后注入延迟，对比\textbf{朴素按时间戳执行}（延迟即错位）与 \textbf{ADG 顺序执行}（保持依赖、谁慢等谁）。复用前文的 \verb@cbs@、\verb@find_first_conflict@。

\begin{codebox}[Python]{ADG 执行：把安全性从绝对时间转成相对顺序}
def build_adg(paths):
    """简化 ADG:记录每个格子按计划被访问的先后次序。
    后访问者进入该格前,须等先访问者离开——这就是跨机顺序依赖。"""
    visits = {}                                   # cell -> [(进入步, 机器人)]
    for i, p in enumerate(paths):
        seen = None
        for t, v in enumerate(p):
            if v != seen: visits.setdefault(v, []).append((t, i)); seen = v
    for v in visits: visits[v].sort()             # 按计划访问时间排序 = 依赖顺序
    return visits

def execute_naive(paths, delay_agent, delay_steps):
    """朴素按时间戳执行:被延迟者整体后移 delay_steps。返回是否无碰撞。"""
    T = max(len(p) for p in paths) + delay_steps + 2; pos = []
    for i, p in enumerate(paths):
        pp = ([p[0]] * delay_steps + list(p)) if i == delay_agent else list(p)
        pos.append(pp + [pp[-1]] * (T - len(pp)))
    return find_first_conflict(pos) is None

def execute_adg(paths, visits, delay_agent, delay_steps):
    """按 ADG 执行:进入格子 v 前须等所有'计划中更早访问 v 的机器人'离开 v。"""
    k = len(paths); comp = []                      # 各机'压缩路径'(去掉原地等待)
    for p in paths:
        c, seen = [], None
        for v in p:
            if v != seen: c.append(v); seen = v
        comp.append(c)
    rank = {}                                      # (格子) -> {机器人: 访问次序}
    for v, lst in visits.items():
        for r, (t, i) in enumerate(lst): rank.setdefault(v, {})[i] = r
    ptr = [0] * k; cur = [comp[i][0] for i in range(k)]
    left = {v: 0 for v in visits}                  # 某格已被多少个'更早访问者'离开
    start = [delay_steps if i == delay_agent else 0 for i in range(k)]
    for step in range(200):
        if all(ptr[i] == len(comp[i]) - 1 for i in range(k)): return True, True
        occ = {cur[i] for i in range(k)}
        for i in range(k):
            if step < start[i] or ptr[i] >= len(comp[i]) - 1: continue
            nxt = comp[i][ptr[i] + 1]
            if left[nxt] < rank[nxt][i]: continue   # ADG 依赖:更早访问者还没离开 nxt
            if nxt in occ: continue                 # 物理顶点占用
            old = cur[i]; cur[i] = nxt; ptr[i] += 1
            occ.discard(old); occ.add(nxt); left[old] = rank[old][i] + 1
        if len(set(cur)) != k: return False, False  # 顶点碰撞
    return True, all(ptr[i] == len(comp[i]) - 1 for i in range(k))

grid3 = np.zeros((3, 3), dtype=int)
s3, g3 = [(1, 0), (1, 2)], [(1, 2), (1, 0)]        # 两机互换
plan = cbs(grid3, s3, g3)[0]; visits = build_adg(plan)
for D in [0, 2, 4]:
    naive = execute_naive(plan, 0, D); adg, done = execute_adg(plan, visits, 0, D)
    print(f"延迟{D}步: 朴素时间戳执行 无碰撞={naive} | ADG执行 无碰撞={adg} 全到达={done}")
\end{codebox}

\textbf{运行要点}：无延迟时两种执行都安全。一旦机器人 $0$ 被延迟（$2$ 步、$4$ 步），\textbf{朴素按时间戳执行就会碰撞}——因为机器人 $1$ 照原时刻推进，而机器人 $0$ 滞后，原本错开的时空对齐被打破。而 \textbf{ADG 执行始终无碰撞且全员到达}（\cref{thm:mapf-adg}）：机器人 $1$ 在要进入某个“机器人 $0$ 计划中更早经过的格子”前，会自动等机器人 $0$ 先离开，无论它慢多少。这把执行的安全性从“脆弱的绝对时间”转成了“鲁棒的相对顺序”。注意这是简化版 ADG（只演示跨机格子顺序）；完整 MAPF-POST/ADG 还要处理边、运动学限制与安全距离。

\begin{pitfall}[概念误区：以为离散 MAPF 计划可以直接在真机上同步执行]
把 CBS/PIBT 算出的“第 $t$ 步谁在哪格”直接当指令发给真机，假设所有机器人严格同步、每步等长。后果：真机一有速度差异或延迟，基于时间戳的无冲突保证立刻失效，发生碰撞。根本原因：离散 MAPF 的安全性依赖“同步 $+$ 单位步长”假设。正确做法：计划与执行解耦——用 MAPF-POST 后处理成考虑运动学的时序调度，或用 ADG 把安全性转成相对顺序（\cref{thm:mapf-adg}），或用 $k$-robust 在规划时预留延迟容忍。绝不裸用时间戳直发真机。
\end{pitfall}

\begin{pitfall}[思维陷阱：用加大时间余量来对抗延迟]
为应付延迟，给每步加大段固定余量、让机器人慢行。后果：吞吐暴跌（全程龟速），且余量是拍脑袋定的——延迟一旦超过余量照样撞，鲁棒性是假的。根本原因：固定时间余量既浪费又无保证；它没有从根本上解决“绝对时间对齐脆弱”的问题。正确做法：用依赖顺序（ADG）而非时间余量来保证安全——ADG 下机器人按需等待（只在依赖未满足时等），既不浪费、又对任意延迟鲁棒；要在规划层建模就用 $k$-robust 给出可证明的延迟容忍。
\end{pitfall}

\begin{pitfall}[编程陷阱：ADG 只建同机顺序、漏掉跨机的格子访问顺序（或反之）]
构建 ADG 时只连“同一机器人前后动作”的边，忘了连“不同机器人对同一格子的访问先后”的边。后果：执行时后到的机器人不会等先到的离开，照样撞。根本原因：ADG 的安全性来自\textbf{两类}依赖——同机顺序与跨机格子顺序；缺一类就不安全。正确做法：两类边都要建。跨机边：对每个格子，按计划访问时间排序，后访问者“进入”依赖于前访问者“离开”（代码里 \verb@left[nxt] < rank[nxt][i]@ 即在检查这条依赖）。
\end{pitfall}

\begin{exercise}
\begin{enumerate}
  \item \textbf{[实现/扩展]} 给 ADG 执行加入\textbf{随机延迟}（每步每个机器人以概率 $p$ 停一拍）。在多组随机延迟下验证 ADG 执行始终无碰撞、全员到达，而朴素时间戳执行的碰撞率随 $p$ 上升。
  \item \textbf{[实现]} 把 ADG 的依赖从“格子访问顺序”扩展到\textbf{边}：若计划中机器人 $i$ 先走过边 $(u,v)$、$j$ 后走同一条边，加一条依赖边。验证在更密集实例上仍无碰撞。
  \item \textbf{[实现/对比，接 $k$-robust]} 实现一个简单的 $k$-robust 检查：验证“任意单个机器人延迟至多 $k$ 步时是否仍无碰撞”。对比一个普通（$0$-robust）计划与一个人工加了缓冲的计划。
  \item \textbf{[思考/设计：CCBS]} 在连续时间几何设定下，两台圆形机器人沿交叉直线运动可能在某个连续时刻“擦碰”。说明为什么离散 MAPF 的“顶点/边冲突”无法刻画这种连续碰撞，以及 CCBS 用“带时刻的动作冲突 $+$ 不安全时间区间约束”如何解决。
  \item \textbf{[综合设计：MAPD/TAPF]} 设计一个仓库 MAPD 系统的执行栈：上层用 Token Passing 做终身任务分配（或目标匿名时用 TAPF/CBM），中层用 LaCAM/PIBT 快速规划路径，底层用 ADG 做对延迟鲁棒的执行。说明三层如何交接、各自输入输出、以及某机器人长期延迟时如何反馈到上层重分配。
\end{enumerate}
\end{exercise}

% ════════════════════════════════════════════════════════════════
\section{方法全景：学习式 MAPF 与选型决策}
\label{sec:mapf-landscape}

\paragraph{动机：方法很多，怎么选才不踩坑。}
前八节给了一大堆方法——匈牙利、CBBA、CBS、ECBS/EECBS、MAPF-LNS、PP、PIBT、LaCAM、CCBS、ADG……工程里两种典型的错都很常见：\textbf{过度工程}——明明几十台机器人、实时性要求高，却非要上最优 CBS，结果跑不动；\textbf{欠工程}——需要质量保证或完备性的场合，图省事用了优先级规划，结果偶发死锁或解很差还不自知。与此同时，这几年又冒出一条新路线——\textbf{学习式 MAPF}，它不再显式搜索，而是让神经网络从数据里学会“怎么走/怎么搜”，在去中心化、实时、不确定的真实场景里有独特优势，但也带来“无保证、难泛化”的新问题。

\subsection{学习式 MAPF：搜索与规则之外的第三条路}
\label{sec:mapf-learning}

前面的方法可归为两类：\textbf{搜索式}（CBS、ECBS/EECBS、LaCAM、CCBS）与\textbf{规则式}（PP、PIBT）。\textbf{学习式}是第三类，有三种典型用法：

\begin{itemize}
  \item \textbf{直接学策略}：把 MAPF 当成多智能体强化学习问题，每个机器人学一个“看局部观测 $\to$ 选下一步动作”的策略，执行时\textbf{去中心化、反应式}地走。代表是 \textbf{PRIMAL}（Sartoretti 等，RA-L 2019）\cite{sartoretti2019primal}——结合强化学习与模仿学习（用专家 MAPF 规划器的演示来训练），学出完全去中心化的策略，实验做到上千智能体；其终身版 PRIMAL2 面向密集仓库；
  \item \textbf{学习引导搜索}：不取代搜索，而是用学习让搜索更快——比如学习 CBS 里“先消解哪个冲突”、或学习 MAPF-LNS 里“破坏哪个邻域”\cite{huang2022learning}，还有学习“对这张地图该用哪个 MAPF 算法”（算法选择）。这类保留了搜索的保证；
  \item \textbf{学习 $+$ 经典混合}：把学习与经典方法的长处结合，如 LNS2$+$RL 用强化学习指导 LNS2 的修复。
\end{itemize}

\textbf{通信学习}这条线常用图神经网络（GNN）让机器人学会“和邻居交换什么信息以协调”，如 MAGAT（消息感知图注意力网络，Li、Lin、Liu、Prorok，RA-L 2021）\cite{li2021magat}——这与共识章的通信图、后续的 MARL 直接接续。学习式的取舍很鲜明：\textbf{优点}是去中心化、亚秒反应、能在线应对不确定与动态、可扩展；\textbf{代价}是\textbf{没有完备性/最优性保证}、训练成本高、且面临\textbf{泛化/分布偏移}。

\begin{table}[H]\centering\small
\caption{搜索式 / 规则式 / 学习式三类范式的互补}
\label{tab:mapf-paradigms}
\begin{tabular}{llllll}
\toprule
范式 & 代表 & 怎么求解 & 完备/最优 & 规模/实时 & 软肋 \\
\midrule
搜索式 & CBS、EECBS、LaCAM、CCBS & 在约束/配置空间搜索 & 强保证 & 中（冲突密时受限） & 约束树可能爆 \\
规则式 & PP、PIBT & 按固定规则单步生成 & 弱/无（PIBT 限双连通） & 极强、亚毫秒 & 短视、不最优 \\
学习式 & PRIMAL、MAGAT & 神经网络学策略/启发 & 无保证 & 去中心化反应 & 训练成本、泛化/分布偏移 \\
\bottomrule
\end{tabular}
\end{table}

正因互补，前沿趋势是\textbf{取长补短的混合}：用学习引导搜索（保留搜索的保证、用学习提速）、或学习 $+$ LNS。理解这张表，你就明白为什么没有“银弹”——选型的本质，就是按你的问题在这三类的长短板之间做匹配。

\subsection{本章方法全景：一张总表}
\label{sec:mapf-mastertable}

\begin{table}[H]\centering\small
\caption{本章方法全景总表（横看类别、纵看取舍）}
\label{tab:mapf-master}
\begin{tabular}{llllll}
\toprule
方法 & 类别 & 最优性 & 完备性 & 规模 & 何时用 \\
\midrule
匈牙利 & 分配·集中 & 最优 & --- & 中（$O(n^3)$） & 一对一一次性分配、要最优 \\
拍卖 & 分配·可分布 & $\varepsilon$-近似 & --- & 大 & 要可分布、近最优 \\
CBBA & 分配·分布 & $50\%$（DMG） & --- & 大 & 无中央节点、多任务序列分配 \\
子模贪婪 & 分配 & $1{-}1/e$ / $\tfrac12$ & --- & 大 & 覆盖/传感等子模效用 \\
联合 A\* & 路径·耦合 & 最优 & 完备 & 极小（$b^k$） & 教学/极少智能体 \\
CBS & 路径·搜索 & 最优 & 完备 & 小-中 & 要最优、冲突不太密 \\
ECBS/EECBS & 路径·搜索 & $\le w\cdot$最优 & 完备 & 中-大 & 要质量保证 $+$ 规模 \\
MAPF-LNS & 路径·anytime & 趋近最优（无界） & --- & 大 & 先可行再改进、随时可停 \\
LNS2 & 路径·修复 & 无界次优 & --- & 极大 & 连可行解都难找时 \\
优先级规划 & 路径·规则 & 次优 & 不完备 & 大 & 不拥挤、要快、可接受次优 \\
PIBT & 路径·规则 & 次优（短视） & 双连通图上完备 & 极大、实时 & 上千台、亚毫秒、终身 \\
LaCAM / LaCAM\textsuperscript{*} & 路径·搜索 & 次优 / anytime最优 & 完备 & 极大 & 要完备 $+$ 规模（\textsuperscript{*}再要最优） \\
CCBS & 路径·连续时间 & 最优 & 完备 & 小-中 & 连续时间、几何体、运动学 \\
学习式（PRIMAL 等） & 路径·学习 & 无保证 & 无保证 & 大、反应式 & 去中心化、在线、不确定环境 \\
\bottomrule
\end{tabular}
\end{table}

\subsection{选型决策框架}
\label{sec:mapf-selection}

拿到一个问题，按这个顺序问下去：

\begin{enumerate}
  \item \textbf{是分配还是路径？}（本章两张图、两类问题）——完整系统两者都要（回 \cref{sec:mapf-submod}）；
  \item \textbf{分配}：一对一一次性？$\to$ 匈牙利（最优）。多任务序列、要分布式？$\to$ CBBA（确认 DMG）。覆盖/传感类子模效用？$\to$ 子模贪婪。需多机协作（MR）？$\to$ 这是另一类（联盟形成）；
  \item \textbf{路径，先问规模与实时性}：几十台、离线、要最优 $\to$ CBS（$+$ICBS/CBSH）。要质量保证又要规模 $\to$ ECBS/EECBS。要“先可行再逼近最优” $\to$ MAPF-LNS / LaCAM\textsuperscript{*}。上千台、实时、终身 $\to$ PIBT/LaCAM；超大规模或连可行解都难 $\to$ LNS2/PIBT；
  \item \textbf{再问场景特性}：连续时间/几何体/运动学 $\to$ CCBS。去中心化、在线、强不确定 $\to$ 学习式（或学习引导的经典法）；
  \item \textbf{最后问执行}：要上真机？$\to$ 计划用 ADG/MAPF-POST 转成鲁棒执行，或规划时上 $k$-robust（回 \cref{sec:mapf-exec}）。终身取送货？$\to$ 上层 MAPD/Token Passing，目标匿名则 TAPF/CBM。
\end{enumerate}

\begin{codebox}[Python]{把选型框架写成一个决策函数}
def recommend_mapf_solver(num_agents, need_optimal, need_guarantee,
                          lifelong=False, continuous_time=False,
                          decentralized=False, real_robot=False):
    """根据问题特征推荐 MAPF 方法。返回(求解器, 执行层)。"""
    if continuous_time:
        solver = "CCBS(连续时间+几何体,SIPP 低层)"
    elif decentralized:
        solver = "学习式(PRIMAL 等)或 PIBT(去中心化、反应式;注意学习式无保证)"
    elif need_optimal:
        solver = "CBS(+ICBS/CBSH 提速)" if num_agents <= 40 else "LaCAM*(anytime 最终最优)"
    elif need_guarantee:                                  # 要有界质量保证
        solver = "ECBS/EECBS(SOC <= w*最优)"
    elif num_agents >= 500 or lifelong:
        solver = "PIBT/LaCAM(上千台、实时、终身);可行性难则 LNS2"
    else:
        solver = "MAPF-LNS(anytime,先可行再改进)或优先级规划(不拥挤时)"
    execution = ("ADG/MAPF-POST 转鲁棒执行" + ("(终身:上层 Token Passing/TAPF-CBM)" if lifelong else "")) \
                if real_robot else "仿真同步执行即可"
    return solver, execution

scenarios = [
    dict(num_agents=8,    need_optimal=True,  need_guarantee=False),
    dict(num_agents=200,  need_optimal=False, need_guarantee=True),
    dict(num_agents=2000, need_optimal=False, need_guarantee=False, lifelong=True, real_robot=True),
    dict(num_agents=12,   need_optimal=True,  need_guarantee=False, continuous_time=True),
    dict(num_agents=100,  need_optimal=False, need_guarantee=False, decentralized=True),
]
for sc in scenarios:
    solver, execu = recommend_mapf_solver(**sc)
    print(f"{sc}\n  -> 求解器: {solver}\n     执行层: {execu}\n")
\end{codebox}

\textbf{运行要点}：这个函数把决策框架变成了可调用的代码——给它问题特征，它返回推荐的求解器与执行层。它不是什么“智能”算法，而是把全章对“每个方法拿什么换什么”的理解结构化下来。

\paragraph{综合选型走查。}
拿一个具体场景走一遍：\textbf{某电商仓库，$800$ 台 AGV，货位网格化，任务（取货$\to$打包）持续到达，要上真机、有通信网络}。按框架：这是“分配 $+$ 路径”完整系统；分配上，任务不断来、无中央调度器最稳，用 \textbf{CBBA}（或终身设定下的 Token Passing）在 AGV 通信图上做；路径上，$800$ 台 $+$ 实时 $+$ 终身，选 \textbf{PIBT/LaCAM} 做快速规划，若某时段特别拥挤、可行性吃紧则用 \textbf{LNS2} 修复；上真机有延迟/速度差，计划用 \textbf{ADG} 转成对延迟鲁棒的执行；整个系统是 \textbf{MAPD} 的终身循环。于是落地方案 $=$ CBBA/TP 分配 $+$ PIBT/LaCAM 路径 $+$ ADG 执行，外加 LNS2 兜底可行性。

\begin{pitfall}[概念误区：以为学习式 MAPF 是“免费午餐”，能直接替代经典方法]
看到 PRIMAL 等能上千智能体、去中心化，就以为可以无脑用它替代 CBS/LaCAM，且期望同样可靠。后果：在训练分布外的地图密度/团队规模上成功率骤降、出现死锁或绕路；且无法给出任何完备性/最优性保证，质量敏感场景翻车。根本原因：学习式用“无保证 $+$ 训练/泛化成本”换“去中心化 $+$ 反应式”。正确做法：把学习式用在它的强项（去中心化反应、引导经典搜索、混合如 LNS2$+$RL），质量/完备性要求高时仍用搜索式；评估时务必测训练分布外的泛化。
\end{pitfall}

\begin{pitfall}[思维陷阱：不按需选型——过度工程或欠工程]
不分场景一律上最优 CBS（过度），或图省事一律用优先级规划（欠）。后果：过度工程在大规模/实时场景跑不动；欠工程在需要保证/完备的场景偶发死锁、解很差。根本原因：没有按“规模 $\times$ 最优性需求 $\times$ 是否终身 $\times$ 真机”匹配方法。正确做法：走 \cref{sec:mapf-selection} 的决策框架（或用 \verb@recommend_mapf_solver@ 起步），按问题特征落点。
\end{pitfall}

\begin{pitfall}[编程陷阱：把全景表的“规模”列当成绝对阈值硬编码]
把“CBS 只能到几十台”“PIBT 上千台”当成精确硬阈值写进代码，不做实测。后果：实际规模上限高度依赖地图结构、冲突密度、时限与实现——硬编码阈值会在你的具体场景上误判。根本原因：表里的“规模”是数量级的经验指引，不是与你的地图/硬件无关的常数。正确做法：用全景表（\cref{tab:mapf-master}）做\textbf{初选}，再在自己的实例分布上跑基准（求解时间、成功率、SOC）确定真实边界与时限回退策略。
\end{pitfall}

\begin{exercise}
\begin{enumerate}
  \item \textbf{[实现/扩展]} 给 \verb@recommend_mapf_solver@ 加更多维度：通信是否受限、是否异构机型、是否有预算约束，再补上分配侧的推荐，做成一个完整的“分配 $+$ 路径 $+$ 执行”选型助手。
  \item \textbf{[分析]} 选三个真实系统（亚马逊仓库、无人机表演、自动化港口），分别用 \cref{tab:mapf-master} 与决策框架推断它们最可能用了哪类方法，并说明理由。
  \item \textbf{[实现/对比]} 复现一个\textbf{最小的学习式 MAPF}：在小栅格上用表格 Q-learning 或简单策略梯度，让 $2$--$3$ 个智能体学“看局部观测选动作”的去中心化策略，和 CBS 对比成功率与路径质量。
  \item \textbf{[思考]} 设计一个“学习 $+$ 搜索”接口：学习模块（如学习 CBS 选哪个冲突、学习 LNS 选哪个邻域）的输入输出是什么、如何保证加入学习后仍不破坏原算法的最优性/有界次优性。
\end{enumerate}
\end{exercise}

% ════════════════════════════════════════════════════════════════
\section{本章常见误解汇总}
\label{sec:mapf-misconceptions}

\begin{table}[H]\centering\small
\caption{本章常见误解与正确理解}
\label{tab:mapf-misconceptions}
\begin{tabular}{p{0.42\linewidth}p{0.52\linewidth}}
\toprule
常见误解 & 正确理解 \\
\midrule
任务分配就是“每台机器人挑最近的任务” & 那是贪婪，有全局耦合时显著次优；一对一一次性分配有最优的匈牙利（$O(n^3)$） \\
所有任务分配问题都能最优求解 & 仅 ST-SR-IA（一对一一次性）多项式最优（LAP）；带 MR、TA 的一般 NP-难 \\
CBBA 是“分布式的匈牙利”，给最优解 & CBBA 等价于集中式\textbf{贪婪}（SGA），保证是 $50\%$ 界（DMG 下），不是最优 \\
CBBA 释放被抢任务时只释放那一个 & 必须连同 bundle 里其后的所有任务一起释放，否则保证失效、易活锁 \\
CBBA 对任何得分函数都有 $50\%$ 界 & 仅当得分满足 DMG（子模）；违反 DMG 时保证失效，需 warping 单调化 \\
多机路径规划 $=$ 每台机器人各跑 A\* & 那必然死锁；MAPF 的难度全在“无碰撞”耦合约束 \\
无碰撞只要“不在同一时刻占同一格” & 还有边/交换冲突（两机在一条边上对向交换）；只查顶点会让对向机器人穿模 \\
sum-of-costs 和 makespan 是一回事 & 前者各路径代价之和（优化吞吐），后者最长路径（优化全部完成时间）；都 NP-难 \\
CBS 发现冲突让一方避让即可 & 必须分裂出“让 $i$”和“让 $j$”两个子节点都展开，否则退化成 PP \\
优先级规划/PIBT 给的是最优解 & 都在解耦侧，用最优性换可扩展；PIBT 还是短视（只一步） \\
优先级规划总能找到解 & 不完备——高优先级机器人停在咽喉要道可堵死低优先级者 \\
PIBT 在任何地图上都完备 & 完备性要求“任意相邻顶点对在某简单环上”（如双连通） \\
子模性就是凸性 & 子模是集合函数的“边际收益递减”，与连续凸性是不同概念 \\
任务分配的图和路径规划的图是同一张 & 两张本质不同的图：通信图（节点$=$机器人）与环境/栅格图（节点$=$位置） \\
分配用直线距离最优，系统就最优 & 真实代价是无冲突路径长度；拥挤时差距巨大，需迭代/联合 \\
次优 MAPF 都没有质量保证 & 有界次优（ECBS/EECBS）保证 SOC $\le w\cdot$ 最优（来自焦点搜索的下界） \\
anytime 算法就是“跑久点的最优算法” & MAPF-LNS 随时可停且总有合法解，通过破坏-修复在已有解上局部改进 \\
离散 MAPF 计划可直接发给真机同步执行 & 真机有延迟/速度差/形状，时间戳对齐会失效；须用 ADG/MAPF-POST/$k$-robust \\
用加大时间余量就能让执行鲁棒 & 既浪费又无保证（延迟超余量照撞）；ADG 用依赖顺序按需等待 \\
任务分配与路径规划必须串行分两步 & 目标匿名时 TAPF/CBM 在一个 CBS 式框架里用最小费用流联合求解 \\
学习式 MAPF 能直接取代经典方法且一样可靠 & 它用“无保证$+$训练/泛化成本”换去中心化反应式；质量要求高仍用搜索式 \\
一个方法可以打天下 & 每个方法的取舍是为特定场景设计的；要按规模$\times$最优性$\times$终身$\times$真机选型 \\
\bottomrule
\end{tabular}
\end{table}

% ════════════════════════════════════════════════════════════════
\section{本章小结}
\label{sec:mapf-summary}

本章把多机协调从连续控制层往上推到\textbf{离散决策层}。

\textbf{任务分配（\cref{sec:mapf-assign,sec:mapf-cbba}）}：谁去做哪件任务？一对一一次性分配是线性指派问题，有最优的匈牙利算法（$O(n^3)$，\cref{thm:mapf-tu}）；但工程上更要可扩展与分布式，于是有拍卖（价格作协调信号）、顺序贪婪（SGA，DMG 下 $50\%$ 界）。Gerkey--Matari\'c 三维分类法告诉你问题属于哪一类。CBBA 把拍卖与共识焊接——本地贪婪构建 bundle，再用带时间戳的最大值共识消解冲突；在 DMG 条件下它等价于集中贪婪 SGA，继承 $50\%$ 界（\cref{thm:mapf-cbba}），收敛轮数正比于通信图直径。

\textbf{路径规划（\cref{sec:mapf-def,sec:mapf-cbs,sec:mapf-scalable}）}：怎样走不撞？MAPF 形式化了顶点冲突与边冲突、sum-of-costs 与 makespan，而最优求解是 NP-难。把算法摆在“耦合$\leftrightarrow$解耦”谱上看最清楚：联合 A\* 最优但 $5^N$ 爆炸；独立 A\* 快但死锁。CBS 居中——低层单体 A\*（解耦的效率）$+$ 高层约束树两子节点穷尽（耦合的最优），拿到最优 $+$ 完备（\cref{thm:mapf-cbs}）。要上千台/实时/终身，就回到解耦侧：PP（快但不完备）、PIBT（优先级继承 $+$ 回溯，亚毫秒/上千台）、LaCAM（惰性后继搜索恢复完备）、LaCAM\textsuperscript{*}（anytime 逼近最优）。

\textbf{边界与统一（\cref{sec:mapf-submod}）}：任务分配的种种贪婪保证（$50\%$、$1-1/e$）背后是同一个\textbf{子模性}；分配与路径耦合成一个完整系统；本章的离散 MAPF 与连续编队是同一“共享空间不相撞”问题在不同抽象层的建模。

\textbf{质量$\times$规模与真机执行（\cref{sec:mapf-bounded,sec:mapf-exec}）}：最优 CBS 与极速 PIBT 之间有最有用的中间地带——\textbf{有界次优}（ECBS/EECBS，焦点搜索保证 SOC $\le w\cdot$ 最优）与 \textbf{anytime}（MAPF-LNS）。把离散计划接到真机，还要跨过“同步、单位步、点机器人”三个理想假设：连续时间用 CCBS（SIPP 安全区间），延迟鲁棒执行用 ADG（\cref{thm:mapf-adg}）或 $k$-robust，终身取送货用 MAPD/Token Passing、目标匿名时用 TAPF/CBM。

回到章首的两个“如果跳过会怎样”：现在你知道贪婪/拍卖分配比最优差一截是\textbf{分布式必须付的税}（$50\%$ 界），也知道各自独立规划必然死锁、必须显式处理顶点 $+$ 边冲突。更重要的是——\textbf{减耦合换可扩展、保耦合换最优，是和共识章“全局协调 vs 局部自主”、ADMM 的 $\rho$ 旋钮同源的同一根权衡}；而子模性则是让“分工”可被高效逼近的统一数学根源。

\subsection*{符号表}
\begin{table}[H]\centering\small
\begin{tabular}{ll}
\toprule
记号 & 含义 \\
\midrule
$N,\ M$ & 机器人数 / 任务数 \\
$\mathbf{C}=[c_{ij}],\ \mathbf{U}=[u_{ij}]$ & 代价矩阵 / 效用矩阵 \\
$x_{ij}\in\{0,1\}$ & 指派变量（$x_{ij}=1$ 表示机器人 $i$ 做任务 $j$） \\
$p_j,\ \varepsilon$ & 拍卖中任务 $j$ 的价格 / 松弛量（$\varepsilon$-最优差距 $\le N\varepsilon$） \\
$\mathbf{b}_i,\ \mathbf{p}_i$ & CBBA 的 bundle（加入序）/ path（执行序） \\
$\mathbf{y}_i,\ \mathbf{z}_i,\ \mathbf{s}_i$ & CBBA 中标价 / 中标者 / 时间戳向量 \\
$c_{ij}[\mathbf{p}_i]$ & 把任务 $j$ 插入 path $\mathbf{p}_i$ 的最佳边际得分 \\
$D$ & 通信图直径（CBBA 收敛轮数 $\propto D$） \\
$\mathcal{G}=(V,E)$ & 环境/栅格图（顶点$=$位置，边$=$相邻移动） \\
$s_i,\ g_i,\ \pi_i$ & 机器人 $i$ 的起点 / 目标 / 路径 \\
$\langle i,j,v,t\rangle$ & 顶点冲突（同时刻同顶点） \\
$\langle i,j,u,v,t\rangle$ & 边/交换冲突（同时刻对穿一条边） \\
$\mathrm{SOC},\ \mathrm{makespan}$ & $\sum_i\mathrm{cost}(\pi_i)$ / $\max_i\mathrm{cost}(\pi_i)$ \\
$w$ & 次优因子（$w\ge 1$，SOC $\le w\cdot$ 最优） \\
$f_{\min}$ & OPEN 表的最小 $f$ 值（焦点搜索的最优下界） \\
$k$ & $k$-robust 的延迟容忍步数 \\
\bottomrule
\end{tabular}
\end{table}

\subsection*{定理速查表}
\begin{table}[H]\centering\small
\begin{tabular}{lll}
\toprule
结论 & 内容 & 出处 \\
\midrule
\cref{thm:mapf-tu} & LAP 全幺模 $\Rightarrow$ 多项式最优（匈牙利 $O(n^3)$） & \cref{eq:mapf-lap} \\
\cref{thm:mapf-auction} & 拍卖总效用与最优差距 $\le N\varepsilon$ & \cref{sec:mapf-auction} \\
\cref{thm:mapf-sga} & DMG 下 SGA $\ge\tfrac12$ 最优 & \cref{def:mapf-dmg} \\
\cref{thm:mapf-cbba} & DMG 下 CBBA $=$ SGA，继承 $50\%$ 界 & \cref{eq:mapf-cbba-half} \\
\cref{thm:mapf-cbs} & CBS 最优且完备（低层最优$+$最优优先$+$两子节点穷尽） & \cref{sec:mapf-cbs} \\
\cref{thm:mapf-pibt} & 相邻顶点对都在简单环上 $\Rightarrow$ PIBT 有限时间到达 & \cref{sec:mapf-pibt} \\
\cref{thm:mapf-submod-bounds} & 基数 $\to 1-1/e$、拟阵 $\to\tfrac12$ & \cref{sec:mapf-submod} \\
\cref{thm:mapf-focal} & 焦点搜索保证 SOC $\le w\cdot$ 最优 & \cref{eq:mapf-focal} \\
\cref{thm:mapf-adg} & ADG 一致顺序执行 $\Rightarrow$ 任意延迟下无碰撞 & \cref{sec:mapf-adg} \\
\bottomrule
\end{tabular}
\end{table}

\subsection*{知识点总表}
\begin{table}[H]\centering\small
\begin{tabular}{lllc}
\toprule
知识点 & 核心要点 & 对应节 & 难度 \\
\midrule
线性指派问题 & 全幺模 $\to$ 多项式最优，匈牙利 $O(n^3)$ & \cref{sec:mapf-assign} & $\star\star\star$ \\
Gerkey--Matari\'c 分类 & ST/MT、SR/MR、IA/TA & \cref{sec:mapf-assign} & $\star\star$ \\
拍卖算法 & 价格作协调信号；$\varepsilon$-近似差距 $\le N\varepsilon$ & \cref{sec:mapf-assign} & $\star\star\star$ \\
顺序贪婪与 DMG & SGA 每轮挑边际最高对；DMG 下 $50\%$ 界 & \cref{sec:mapf-assign} & $\star\star\star$ \\
CBBA 两阶段 & bundle 构建（本地贪婪）$+$ 共识冲突消解 & \cref{sec:mapf-cbba} & $\star\star\star\star$ \\
CBBA $=$ SGA 与 $50\%$ 界 & DMG 下等价集中贪婪 & \cref{sec:mapf-cbba} & $\star\star\star\star$ \\
按 bundle 顺序释放 & 被抢任务连其后全释放；保出价有效性 & \cref{sec:mapf-cbba} & $\star\star\star\star$ \\
CBBA 收敛 & 正比通信图直径 $D$ & \cref{sec:mapf-cbba} & $\star\star\star$ \\
MAPF 形式化 & 图$+$机器人$+$起讫$+$离散步 & \cref{sec:mapf-def} & $\star\star$ \\
顶点冲突与边冲突 & 漏边冲突 $\to$ 穿模 & \cref{sec:mapf-def} & $\star\star\star$ \\
SOC vs makespan & 总和 vs 最长；都 NP-难 & \cref{sec:mapf-def} & $\star\star$ \\
耦合$\leftrightarrow$解耦谱 & 联合 A\* 爆 $\leftrightarrow$ 独立 A\* 死锁；CBS 居中 & \cref{sec:mapf-def} & $\star\star\star$ \\
CBS 两层结构 & 高层约束树 $+$ 低层单体时空 A\* & \cref{sec:mapf-cbs} & $\star\star\star\star$ \\
CBS 最优性 & 低层最优$+$最优优先$+$两子节点穷尽 & \cref{sec:mapf-cbs} & $\star\star\star\star$ \\
优先级规划$+$预留表 & 低让高；快但不完备/次优 & \cref{sec:mapf-scalable} & $\star\star\star$ \\
PIBT & 单步配置生成；优先级继承$+$回溯；亚毫秒 & \cref{sec:mapf-scalable} & $\star\star\star\star$ \\
PIBT 完备性 & 相邻顶点对在简单环上（双连通） & \cref{sec:mapf-scalable} & $\star\star\star\star$ \\
LaCAM / LaCAM\textsuperscript{*} & 惰性后继搜索完备；\textsuperscript{*} anytime 最终最优 & \cref{sec:mapf-scalable} & $\star\star\star\star$ \\
子模性与近似界 & 基数 $\to 1-1/e$、拟阵 $\to\tfrac12$ & \cref{sec:mapf-submod} & $\star\star\star\star$ \\
焦点搜索 & FOCAL$=\{f\le w f_{\min}\}$；保 $w$-界 & \cref{sec:mapf-bounded} & $\star\star\star\star$ \\
ECBS / EECBS & 高低层焦点搜索 / 高层显式估计搜索$+$在线学习 & \cref{sec:mapf-bounded} & $\star\star\star\star$ \\
MAPF-LNS / LNS2 & 破坏-修复 anytime / 修复可行性 & \cref{sec:mapf-bounded} & $\star\star\star\star$ \\
连续时间 MAPF / CCBS & MAPF\textsubscript{R}；CBS$+$SIPP$+$连续冲突 & \cref{sec:mapf-exec} & $\star\star\star\star$ \\
SIPP & (顶点, 安全区间)状态空间 & \cref{sec:mapf-exec} & $\star\star\star$ \\
动作依赖图 ADG & 同机$+$跨机依赖；延迟下无碰撞 & \cref{sec:mapf-exec} & $\star\star\star\star$ \\
$k$-robust MAPF & 规划时预留缓冲，容忍 $k$ 步延迟 & \cref{sec:mapf-exec} & $\star\star\star$ \\
MAPD / TAPF-CBM & 终身取送货 TP；目标匿名时 CBM 联合 & \cref{sec:mapf-exec} & $\star\star\star\star$ \\
学习式 MAPF & PRIMAL、GNN 通信、学习引导搜索；无保证 & \cref{sec:mapf-landscape} & $\star\star\star$ \\
方法全景与选型 & 搜索/规则/学习三类；按四维选型 & \cref{sec:mapf-landscape} & $\star\star\star$ \\
\bottomrule
\end{tabular}
\end{table}

% ════════════════════════════════════════════════════════════════
\section{延伸阅读}
\label{sec:mapf-further}

\textbf{教材（打地基）}：LaValle 的《Planning Algorithms》\cite{lavalle2006planning}（运动规划与搜索的权威教材，A\*、组合规划的基础都在此）。

\textbf{综述/经典论文（看脉络）}：Gerkey \& Matari\'c 的 MRTA 分类法\cite{gerkey2004taxonomy}（分类法源头，做任务分配必读）；Choi、Brunet、How 的 CBBA 原文\cite{choi2009cbba}（含 DMG、$50\%$ 界、收敛证明）；Sharon 等人的 CBS 原文\cite{sharon2015cbs}（最优 MAPF 标准算法）；Stern 等人的 MAPF 定义与基准综述\cite{stern2019mapf}（公共语言）；Yu \& LaValle 的 NP-难证明\cite{yu2013structure}；Nemhauser、Wolsey、Fisher 的子模贪婪 $1-1/e$ 界\cite{nemhauser1978analysis}。

\textbf{前沿（看标杆）}：Okumura 的 LaCAM\cite{okumura2023lacam} 与 LaCAM\textsuperscript{*}\cite{okumura2023lacamstar}；Okumura 等人的 PIBT\cite{okumura2022pibt}（$2023$ MAPF 竞赛冠军核心）；Xu、Garimella、Tzoumas 的分布式抗失效子模任务分配\cite{xu2025distributed}\pz{存疑：该文实为“通信/计算高效”的分布式子模（RAG），非“抗失效”——详见 \cref{sec:mapf-submod-theory} 处批注。} 与 Corah \& Michael 的分布式子模\cite{corah2019distributed}；Li、Ruml、Koenig 的 EECBS\cite{li2021eecbs}；Li 等人的 MAPF-LNS\cite{li2021lns} 与 LNS2\cite{li2022lns2}；Andreychuk 等人的 CCBS\cite{andreychuk2022ccbs} 与 Phillips \& Likhachev 的 SIPP\cite{phillips2011sipp}；H\"onig 等人的 MAPF-POST\cite{honig2016mapfpost} 与 ADG\cite{honig2019persistent}；Ma 等人的 MAPD/Token Passing\cite{ma2017lifelong} 与 Ma \& Koenig 的 TAPF/CBM\cite{ma2016optimal}；Sartoretti 等人的 PRIMAL\cite{sartoretti2019primal} 与 Li 等人的 MAGAT\cite{li2021magat}。

\textbf{工程（动手查）}：\verb@scipy.optimize.linear_sum_assignment@（匈牙利实现）；\verb@pypibt@（Okumura 本人的 PIBT 最小 Python 实现）；EECBS 与 MAPF-LNS2 的官方 C++ 实现；MAPF benchmark（movingai.com / Stern 等 $2019$）；\verb@networkx@（构造与分析通信图）。

% ════════════════════════════════════════════════════════════════
\section{本章与后续章节的关系}
\label{sec:mapf-next}

\begin{table}[H]\centering\small
\begin{tabular}{p{0.20\linewidth}p{0.46\linewidth}p{0.28\linewidth}}
\toprule
后续主题 & 关系 & 本章铺垫的知识点 \\
\midrule
分布式 MPC 编队 & MAPF 出离散路点，MPC 平滑跟踪；CCBS 连续时间视角过渡到连续控制层 & MAPF 路点、CCBS、耦合$\leftrightarrow$解耦权衡 \\
协同搬运与力控 & MR（多机协作）任务 $=$ 超可加/违反子模，正是本章子模框架的边界 & Gerkey MR 类、子模性边界 \\
异构多机协同 & 任务分配要匹配异构机型能力——把“能力”编码进 CBBA 的得分函数 & CBBA、Gerkey 分类、得分函数接口 \\
MARL（\cref{ch:plan-safemarl} 起） & 学习式的任务分配/路径协调，用策略替代显式搜索 & MAPF、分配作为可学习决策 \\
综合实战 & 仓库调度复用 CBBA$+$CBS$+$LNS 与 ADG 执行，扩成终身分配$\leftrightarrow$路径$\leftrightarrow$执行环 & 全章，尤其 MAPD、ADG \\
\bottomrule
\end{tabular}
\end{table}

\section*{故障排查手册}
\label{sec:mapf-troubleshoot}

\begin{table}[H]\centering\small
\begin{tabular}{p{0.22\linewidth}p{0.30\linewidth}p{0.34\linewidth}p{0.08\linewidth}}
\toprule
症状 & 可能原因 & 排查步骤 & 相关节 \\
\midrule
CBBA 不收敛/反复抖动 & 得分非 DMG；或释放没按 bundle 顺序 & 验证得分边际非增；检查是否连同其后任务一起释放；加 warping & \cref{sec:mapf-cbba} \\
CBBA 收敛了但仍有冲突 & 共识/时间戳规则错；或通信图不连通 & 算通信图 $\lambda_2$ 查连通；检查中标价更新与释放逻辑 & \cref{sec:mapf-cbba} \\
分配总代价远高于预期 & 用了贪婪却期望最优；或代价矩阵行列方向反 & 与匈牙利对比；检查“行$=$机器人、列$=$任务” & \cref{sec:mapf-assign} \\
MAPF 路径执行时相撞 & 漏检边/交换冲突；或目标后占用没处理 & 检查冲突检测的 $t{+}1$ 分支；加目标到达后永久占用 & \cref{sec:mapf-def} \\
CBS 不返回/极慢 & 冲突太多约束树爆；或低层没遵守全部约束 & 确认低层用 CT 节点完整约束集；上 ICBS/CBSH 或改 ECBS & \cref{sec:mapf-cbs} \\
CBS 返回的解不是最优 & 高层没按 cost 最优优先；或只展开一个子节点 & 确认高层按 SOC 出队；确认每个冲突分裂两个子节点 & \cref{sec:mapf-cbs} \\
优先级规划无解 & 高优先级停在咽喉要道堵死；或顺序差 & 换优先级顺序/随机重启；改用 PIBT 或 LaCAM & \cref{sec:mapf-scalable} \\
PIBT 某些机器人到不了目标 & 图非双连通（有死胡同） & 检查“相邻顶点对是否都在简单环上”；加边使图双连通 & \cref{sec:mapf-scalable} \\
终身 MAPF 活锁 & 滑动窗口太短；或优先级反复横跳 & 加长窗口；保证优先级单调；对久停机器人提权 & \cref{sec:mapf-scalable} \\
分配纸面最优、执行却很堵 & 分配用直线距离忽略拥堵 & 改用考虑拥堵的真实路径代价；路径反馈给分配迭代 & \cref{sec:mapf-submod} \\
行为整体错乱 & 把通信图当环境图（或反之） & 代码里分开命名 \verb@comm_graph@/\verb@env_grid@ & \cref{sec:mapf-submod} \\
MAPF-LNS 不改进/极慢 & 邻域纯随机；或修复未把他人当完整障碍 & 改有信息邻域；修复时登记顶点$+$边$+$目标占用并复核 & \cref{sec:mapf-bounded} \\
有界次优解质量没保证 & 把无界放松当有界；或 FOCAL 阈值/下界算错 & 确认按 $f\le w f_{\min}$ 框 FOCAL；确认 $f_{\min}$ 取自 OPEN 真实最小 & \cref{sec:mapf-bounded} \\
离散计划在真机上偶发相撞 & 裸用时间戳同步执行 & 改用 ADG（相对顺序）；或 MAPF-POST；或规划时上 $k$-robust & \cref{sec:mapf-exec} \\
ADG 执行仍碰撞 & 只建同机顺序、漏跨机格子顺序（或漏边依赖） & 对每格按访问序加“后到等先到离开”依赖；补边依赖 & \cref{sec:mapf-exec} \\
学习式 MAPF 某些地图骤然失效 & 测试地图密度/团队规模在训练分布外 & 在目标分布上重训/微调；质量敏感场景回退搜索式 & \cref{sec:mapf-landscape} \\
选了方法却跑不动/解很差 & 过度工程或欠工程 & 走选型框架重新选型；在自己实例上测真实规模边界 & \cref{sec:mapf-landscape} \\
\bottomrule
\end{tabular}
\end{table}

\begin{note}
任务分好了、路也规划好了——但本章给的路径是\textbf{离散}的（每步走一格）、是\textbf{运动学}层面的（没考虑机器人的动力学、惯性、力）。真实机器人不能瞬间转向、急停，编队飞行要平滑、协同搬运要控力。把本章离散决策层的输出，接到能处理动力学与物理交互的\textbf{连续控制层}，正是后续协同运动控制章节的任务——而它跟踪的参考路点，正可以来自本章的 MAPF。离散与连续，在那里交汇。
\end{note}
